\documentclass[11pt,a4paper]{article}
\usepackage[T1]{fontenc}
\usepackage[utf8]{inputenc}
\usepackage[french]{babel}
\usepackage{lmodern}
\usepackage{amsmath,amssymb}
\usepackage[margin=2.35cm]{geometry}
\usepackage{booktabs,longtable,tabularx}
\usepackage{array}
\usepackage{enumitem}
\usepackage{graphicx}
\usepackage{xcolor}
\usepackage{microtype}
\usepackage{float}
\usepackage{hyperref}
% Caractères Unicode présents dans les fragments générés (noms de fenêtres).
\DeclareUnicodeCharacter{2264}{\ensuremath{\le}}
\DeclareUnicodeCharacter{2265}{\ensuremath{\ge}}
\DeclareUnicodeCharacter{03C3}{\ensuremath{\sigma}}
\DeclareUnicodeCharacter{03B3}{\ensuremath{\gamma}}
\DeclareUnicodeCharacter{03C6}{\ensuremath{\varphi}}

\hypersetup{colorlinks=true,linkcolor=blue!45!black,urlcolor=blue!45!black}
\setlist{nosep}
\newcommand{\code}[1]{\texttt{#1}}
\newcommand{\Kaut}{K^{*}_{\mathrm{aut}}}
\newcommand{\fait}{\textbf{[Fait observé]}}
\newcommand{\inference}{\textbf{[Inférence]}}
\newcommand{\hyp}{\textbf{[Hypothèse]}}
\newcommand{\incertitude}{\textbf{[Incertitude]}}

\title{M4.3Live --- rapport de résultats\\
\large Effet rebond : réponse de la production agrégée à une hausse
locale ou globale de la puissance d'extraction}
\author{Programme M4.3Live --- dossier \code{m4\_3live\_credit\_soc}}
\date{17 août 2026}

\begin{document}
\maketitle

\begin{abstract}
L'objectif de cette itération est d'\emph{observer} un éventuel effet
rebond, pas de l'expliquer (prompt §1). Le verdict n'est pas un mot : il
dépend de la portée de l'intervention \emph{et} de l'horizon
d'observation, et les deux dimensions donnent des réponses opposées.

\textbf{À portée partielle et court horizon, l'effet rebond est
observé} : la production agrégée répond 2,3 à 2,9 fois plus fort que la
réponse strictement proportionnelle, avec une significativité de 3 à
48$\sigma$ selon l'intensité, contre un bras de contrôle qui consomme
exactement le même flux aléatoire. \textbf{À portée globale et régime établi, l'effet est
inverse} : une hausse uniforme de 50\,\% du coefficient $A$ ne produit que
$+36{,}1 \pm 0{,}4$\,\% de production agrégée, soit $0{,}72$ fois la
réponse proportionnelle --- et la décomposition montre pourquoi la
somme est sous-proportionnelle alors que chaque entité, elle,
sur-répond : la production par entité est multipliée par $1{,}80$ (pour
un choc de $1{,}50$) mais la population se contracte de 25\,\%.
\textbf{À portée partielle et long horizon, la question n'est pas
tranchable} : la cohorte traitée s'éteint, l'intensité de traitement tend
vers zéro et le dénominateur de toute élasticité disparaît.

\textbf{Mais le signe du résultat global tient à une convention}, et une
ablation dédiée le montre (§\ref{sec:ablation}) : la contraction de
population n'est pas un effet de la technologie, c'est un effet d'échelle du
capital de naissance $K_0$, resté fixe pendant que le système grandissait.
Compenser $K_0$ l'annule entièrement --- population $\times1{,}005$ --- et
la réponse devient alors fortement super-proportionnelle,
$\varepsilon = 2{,}02$. L'explication concurrente, un service d'intérêts
alourdi, est \emph{réfutée} : rapporté à la production, le service baisse.
\textbf{Cette élasticité suit une loi} (§\ref{sec:loi}) :
$\varepsilon = 1/(1-\gamma)$, l'exposant de l'échelle autarcique, vérifié à
mieux que 1,3~\% en $\gamma = 0{,}4$, $0{,}5$ et $0{,}6$ --- et
\emph{démontrée} : le modèle est covariant d'échelle, une hausse de $A$
accompagnée de $K_0 \to c\,K_0$ avec $c = (A'/A)^{1/(1-\gamma)}$ redonne la
même simulation à l'échelle $c$ près, ce qui se vérifie à la précision
machine.

\textbf{Une grandeur unique, la « tension » $T = \Kaut/K_{\text{eq}}$,
résume ces effets d'échelle} (§\ref{sec:tension}) : le rapport entre le
capital d'équilibre d'une entité isolée et le capital auquel le système
tourne réellement. Le système observé tourne à un treizième de son échelle
autarcique, et la mortalité y suit $T^{0{,}67}$ à $1{,}4$\,\% près sur une
plage $\times15$ atteinte par trois leviers différents --- deux
interventions sans rapport ($K_0/4$ et $A\times2$) qui amènent à la même
tension donnent la même mortalité à $0{,}5$\,\% près. \textbf{Mais la tension
n'est pas un paramètre d'état} : un balayage à quatre leviers la met en
échec sur la dépréciation $\delta$, où elle varie d'un facteur 13 pour une
mortalité qui ne bouge que de 14\,\%. Ce qui aligne les quatre leviers est
l'intensité de rotation du crédit, avec $R^2 = 0{,}998$ --- une variable
endogène, qui désigne le canal sans le démontrer.

Un résultat méthodologique conditionne les précédents : le raisonnement du
prompt (§2) selon lequel la portée \code{fraction} « fixe l'intensité de
traitement » ne tient pas dans ce modèle. La part de production détenue
par la cohorte traitée passe de $0{,}198$ à $0{,}022$ en 2000 pas. Les
deux portées dérivent ; \code{fraction} dérive vers le bas.
\end{abstract}

\tableofcontents

% Section partagée par les deux rapports M4.3Live : le modèle et les
% notations, pour que chaque document se lise sans accès au code ni aux
% rapports des lignées antérieures.
\section{Le modèle, et les notations employées}
\label{sec:modele}

\subsection{Ce que simule M4.3Live}

Le modèle décrit une population d'entités économiques abstraites qui
échangent une ressource unique, le \emph{joule}, assimilée ici à un capital.
Il n'y a ni monnaie, ni banque centrale, ni prix : seulement des capitaux,
une fonction de production, et des prêts bilatéraux perpétuels.

Une entité $i$ est entièrement décrite par trois nombres : son capital
$K_i \ge 0$ et sa \emph{technologie} $(A_i, \gamma_i)$, qui fixe la quantité
qu'elle produit à chaque pas de temps,
\begin{equation}
\text{production de } i \;=\; A_i\,K_i^{\gamma_i},
\qquad A_i > 0,\quad 0 < \gamma_i < 1 .
\label{eq:production}
\end{equation}
L'exposant $\gamma_i < 1$ rend les rendements \emph{marginaux} décroissants :
un joule supplémentaire rapporte d'autant moins que le capital est déjà
grand ($\partial^2/\partial K^2 < 0$). La production, elle, continue de
croître avec le capital --- doubler le capital augmente la production, mais
de moins du double. C'est cette concavité qui crée un intérêt à déplacer du
capital d'une entité riche vers une entité pauvre --- le fondement même du
marché du crédit du modèle.

\subsection{Un pas de temps}

Chaque pas déroule huit phases, dans cet ordre invariable :

\begin{enumerate}
\item \textbf{Interventions} --- les changements de paramètres demandés en
  direct sont appliqués (propre à M4.3Live ; voir §\ref{sec:portees} pour
  leurs portées).
\item \textbf{Naissances} --- un nombre $\mathrm{Poisson}(\lambda)$ d'entités
  neuves apparaissent, chacune dotée du capital de naissance $K_0$.
\item \textbf{Choc} --- chaque capital est multiplié par $e^{\xi}$ avec
  $\xi \sim \mathcal{N}(-\sigma^2/2,\ \sigma^2)$, un choc multiplicatif
  d'espérance 1 (il ne crée ni ne détruit de richesse en moyenne).
  \emph{Ordre de grandeur.} L'écart typique se lit sur
  $\mathbb{E}|\xi| = \sigma\sqrt{2/\pi}$ à la correction de moyenne près
  ($-\sigma^2/2$, négligeable ici) : à $\sigma = 0{,}01$, la valeur employée
  dans tout ce programme, $\mathbb{E}|\xi| \approx 8{,}0\times10^{-3}$, donc
  un pas multiplie un capital par $1 \pm 0{,}8\,\%$ en moyenne, et par
  $1 \pm 2{,}0\,\%$ dans 5\,\% des cas. Le choc est un bruit de fond, pas un
  moteur : sur les 2000 pas d'un amorçage, sa contribution \emph{nette}
  cumulée (\code{shock\_gain}) vaut $-0{,}41\,\%$ du capital agrégé final,
  et aucun pas isolé ne déplace \code{K\_tot} de plus de $0{,}16\,\%$.
\item \textbf{Production} --- chaque entité produit selon
  \eqref{eq:production} et \emph{ajoute intégralement} sa production à son
  propre capital.
\item \textbf{Service des intérêts} --- chaque emprunteuse doit
  $\mathrm{due} = \sum_k q_k r_k$, somme sur ses contrats du principal $q_k$
  fois le taux $r_k$. Si son capital ne suffit pas, elle paie ce qu'elle
  peut, tombe à zéro et est marquée en \emph{défaut de liquidité}.
\item \textbf{Dépréciation} --- chaque capital est multiplié par $1-\delta$.
\item \textbf{Marché du crédit} --- $\lfloor \eta(N) \rfloor$ tirages de
  paires, où $N$ est le nombre d'entités éligibles et
  $\eta_{\rho,\beta}(N) = \rho\,N_{\text{ref}}(N/N_{\text{ref}})^{\beta}$
  (la baseline $\rho=1$, $\beta=1$ donne simplement $\eta(N)=N$). Dans chaque
  paire, la plus riche est la \emph{prêteuse}, la plus pauvre
  l'\emph{emprunteuse} ; le prêt ne va jamais dans l'autre sens. Un contrat
  transfère un principal $q$ et fixe un taux $r$, tous deux gelés à vie.
\item \textbf{Faillites en cascade} --- toute entité dont la valeur nette
  $K + \text{créances} - \text{dettes}$ devient négative, ou qui est en
  défaut de liquidité, meurt : ses dettes sont annulées (perte sèche pour
  ses prêteuses, qui peuvent tomber à leur tour) et son capital est détruit.
  On itère jusqu'à point fixe.
\end{enumerate}

Un identifiant d'entité n'est jamais réutilisé ; une entité morte le reste.

\subsection{Notations}

\begin{center}\footnotesize
\begin{tabular}{@{}llp{8.6cm}@{}}
\toprule
\textbf{Symbole} & \textbf{Nom} & \textbf{Définition} \\
\midrule
$K_i$ & capital & Ressource détenue par l'entité $i$. \\
$A_i$ & coefficient d'extraction & Facteur multiplicatif de la production
\eqref{eq:production}. \textbf{Levier primaire} de ce programme. \\
$\gamma_i$ & exposant d'extraction & Exposant de la production : il fixe le
rendement d'échelle. C'est lui qui rend les rendements \emph{marginaux}
décroissants ($0<\gamma<1$), et non le niveau de la production, qui croît
toujours avec le capital. Levier secondaire de ce programme. \\
$(A_i,\gamma_i)$ & technologie & Couple fixé à la naissance, jamais modifié
sauf par une intervention. Reçoit un identifiant entier interne. \\
$\lambda$ & taux de naissance & Espérance du nombre d'entités créées par pas. \\
$K_0$ & capital de naissance & Capital dont une entité neuve est dotée. \\
$\delta$ & dépréciation & Fraction du capital perdue à chaque pas. \\
$\sigma$ & amplitude du choc & Écart-type du choc multiplicatif log-normal. \\
$N$ & taille du bassin & Nombre d'entités éligibles au marché à ce pas. \\
$x,\ y$ & capitaux de la paire & $x = K_b$ (emprunteuse), $y = K_\ell$
(prêteuse), avec $y > x$ par la règle de marché. \\
$C$ & somme conservée & $C = x+y$ ; le transfert la laisse inchangée. \\
$q$ & principal & Montant transféré par le contrat. \\
$r$ & taux & Service versé par pas, perpétuel : l'emprunteuse doit $q\,r$ à
chaque pas, indéfiniment. \\
$\delta^{*}$ & transfert optimal & Principal qui maximise la production
jointe de la paire (§\ref{sec:modele}, éq.~\ref{eq:transfert}). \\
$\lambda^{*}$ & part optimale & Part de $C$ revenant à l'emprunteuse à
l'optimum ; $h(C) = C\lambda^{*}$ est son capital optimal. \\
$\Kaut$ & échelle autarcique & $[A(1-\delta)/\delta]^{1/(1-\gamma)}$ : capital
auquel production et dépréciation s'équilibrent sans marché. Repère
d'échelle du système. \\
\code{prod\_tot} & production agrégée & Somme des productions du pas.
\textbf{Variable d'intérêt de ce programme.} \\
\code{K\_tot}, \code{pop} & agrégats & Capital total et effectif vivant. \\
\code{deaths} & morts & Entités mortes au pas (racines + victimes de cascade). \\
\code{defaults} & défauts & Entités en défaut de liquidité au pas. \\
\code{roots\_insolvency} & racines d'insolvabilité & Entités dont la valeur
nette est passée sous zéro par leurs propres pertes (à distinguer des
victimes de cascade). \\
\code{loan\_volume} & volume de prêt & Somme des principaux transférés au pas. \\
\code{mkt\_rounds} & tirages & Nombre de paires tirées au pas. \\
\code{mkt\_blocked\_dir} & paires refusées & Paires où $\delta^{*} \le 0$ :
l'optimum voudrait faire circuler le capital du pauvre vers le riche, ce que
le sens du prêt interdit. \textbf{Compteur propre à M4.3Live.} \\
\bottomrule
\end{tabular}
\end{center}

\subsection{L'institution de principal de M4.3Live}

C'est le seul point où M4.3Live diffère de son prédécesseur M4.3. Là où M4.3
transférait la moitié de l'écart de capital, $q = (y-x)/2$, M4.3Live
transfère le montant qui maximise la production \emph{jointe} de la paire
immédiatement après l'échange :
\begin{equation}
\delta^{*} \;=\; \arg\max_{-x \le \delta \le y}\;
A_b\,(x+\delta)^{\gamma_b} + A_\ell\,(y-\delta)^{\gamma_\ell}
\;=\; h(C) - x,
\qquad h(C) = C\,\lambda^{*}.
\label{eq:transfert}
\end{equation}
Trois conséquences importent pour lire les deux rapports :

\begin{enumerate}
\item Quand les deux entités partagent la \emph{même} technologie,
  $\lambda^{*} = 1/2$ et $\delta^{*} = (y-x)/2$ exactement : la baseline
  homogène de M4.3Live \emph{est} celle de M4.3.
\item Dès que les technologies diffèrent, $\lambda^{*} \ne 1/2$ et le partage
  est pondéré en faveur de celle dont le \emph{rendement marginal est le plus
  élevé au capital considéré}. Il n'existe pas de « meilleure »
  technologie dans l'absolu : à exposants inégaux, l'ordre des rendements
  marginaux $\gamma A K^{\gamma-1}$ s'inverse avec le capital. Un $\gamma$
  élevé l'emporte aux grands capitaux, un $A$ élevé aux petits, et il existe
  toujours un capital de croisement. C'est pourquoi $\lambda^{*}$ dépend de
  $C$ et non seulement des paramètres.
\item $\delta^{*}$ peut être \emph{négatif}. Comme $\delta^{*} = C\lambda^{*} - x$,
  on a exactement
  \[ \delta^{*} \le 0 \iff \frac{x}{C} \ge \lambda^{*} , \]
  c'est-à-dire : dès que la prêteuse est celle dont le rendement marginal est
  le plus élevé, l'optimum voudrait lui envoyer du capital, et la paire ne
  traite pas. C'est le compteur \code{mkt\_blocked\_dir}, et c'est un canal à
  part entière, pas un détail numérique.
\end{enumerate}

\paragraph{Surplus coopératif.} Le \textbf{surplus coopératif} d'une paire
est le gain de production \emph{jointe} que l'échange produit, à capital
total inchangé :
\begin{equation}
S(x,y) \;=\;
\underbrace{A_b (x+q)^{\gamma_b} + A_\ell (y-q)^{\gamma_\ell}}
          _{\text{après le transfert}}
\;-\;
\underbrace{A_b\,x^{\gamma_b} + A_\ell\,y^{\gamma_\ell}}_{\text{avant}}
\;\ge\; 0 ,
\label{eq:surplus}
\end{equation}
où $q$ est le principal effectivement transféré --- égal à $\delta^{*}$ sauf
quand le plafond mord. La positivité est immédiate quand $q = \delta^{*}$ :
$\delta^{*}$ maximise le premier terme et $\delta = 0$ est admissible ; sous
plafond elle reste vraie car $q$ est entre $0$ et $\delta^{*}$ et le premier
terme est concave en $\delta$. $S$ est nul exactement quand la paire est
déjà à l'optimum ($x/C = \lambda^{*}$), et les paires refusées ne
contribuent pas. La colonne \code{mkt\_surplus} des séries est la somme de
$S$ sur les paires traitées au pas. C'est un \emph{flux de
production par pas}, homogène à \code{prod\_tot}, et non un stock de
capital : le marché ne crée aucun joule, il déplace du capital vers l'endroit
où il produit davantage.

\subsection{Les trois portées d'intervention}
\label{sec:portees}

Une intervention change un paramètre \emph{pendant} qu'une simulation
tourne. Sa \textbf{portée} dit qui elle touche :

\begin{description}
\item[\code{all} (toutes)] --- rétroactif \emph{et} prospectif : toutes les
  entités vivantes reçoivent la nouvelle valeur, et la valeur par défaut à la
  naissance change aussi.
\item[\code{new} (nouvelles)] --- \emph{vintage} technologique : seule la
  valeur par défaut à la naissance change. Les entités déjà vivantes gardent
  la leur pour toujours ; la population se convertit par renouvellement.
\item[\code{fraction} (une fraction $\varphi$)] --- une fraction $\varphi$
  des entités vivantes est tirée au sort une fois, et elle seule est
  modifiée. \textbf{La valeur par défaut à la naissance n'est pas touchée} :
  aucune naissance ne vient reconstituer la cohorte traitée.
\end{description}

Certaines combinaisons n'ont pas de sens et sont refusées explicitement :
$K_0$ n'existe qu'au moment de la naissance, donc \code{all} et \code{new} y
sont le même objet et \code{fraction} n'a pas d'objet ; $\lambda$, $\delta$,
$\sigma$, $\rho$ sont des paramètres de population et non des attributs
d'entité, donc \code{new} y est un alias explicite de \code{all}.


\section{Ce qui est mesuré}

La puissance d'extraction d'une entité est portée par $A$ et $\gamma$ dans
$A\,K^{\gamma}$. $A$ est le levier primaire : il multiplie linéairement la
production à capital donné, quel que soit ce capital. $\gamma$ est le levier
secondaire, et son effet dépend du capital : passer de $\gamma$ à
$\gamma'$ multiplie la production par $K^{\gamma'-\gamma}$. \fait{} Dans le
régime observé, $K/\text{entité} \approx 773$, donc l'augmenter augmente
bien la production, mais d'un facteur qui n'est pas choisi : il faut le
mesurer.

\paragraph{Trois quantités, et un exemple pour les fixer.} Tout ce rapport
tourne autour de trois nombres. On les définit ici une fois, puis on les
illustre immédiatement sur le bras principal de la campagne.

\begin{description}
\item[$h$, l'\emph{horizon}] --- le nombre de pas écoulés depuis
  l'intervention. $h=1$ est le pas où l'intervention agit pour la première
  fois, $h=200$ est deux cents pas plus tard. Toutes les courbes de ce
  rapport sont en $h$, jamais en $t$, pour que les bras se superposent.
\item[$m$, l'\emph{amplitude}] --- de combien la production d'une entité
  \emph{traitée} est multipliée, à capital inchangé :
  $m = A' K^{\gamma'} / A K^{\gamma}$. Pour un levier sur $A$, $m = A'/A$
  exactement. Pour un levier sur $\gamma$, $m = K^{\gamma'-\gamma}$ dépend
  du capital et doit être mesuré.
\item[$p$, la \emph{part traitée}] --- la fraction de la production agrégée
  que les entités traitées auraient produite \emph{sans} traitement (part
  \og ex ante \fg{}).
\end{description}

\noindent\emph{Exemple.} Bras \code{frac\_A150\_phi20} : on tire 20\,\% des
entités vivantes et on multiplie leur $A$ par $1{,}5$, donc $m = 1{,}5$. Ces
entités produisaient $p = 19{,}84\,\%$ de la production totale. Si rien
d'autre ne bougeait dans le système, la production agrégée augmenterait de
$(m-1)\,p = 0{,}5 \times 0{,}1984 = 9{,}92\,\%$ --- c'est ce qu'on appelle la
réponse \textbf{strictement proportionnelle}. Elle sert de mètre étalon :

\begin{equation}
E(h) \;=\; \frac{\big[\mathrm{prod}_{\text{traité}}(h)-\mathrm{prod}_{\text{réf}}(h)\big]
/\;\mathrm{prod}_{\text{réf}}(h)}{(m-1)\,p(h)}
\;=\;
\frac{\text{réponse observée à l'horizon } h}{\text{réponse strictement proportionnelle}},
\label{eq:elasticite}
\end{equation}
avec $p(h)$ remesuré à chaque pas, car la cohorte traitée change de poids au
fil du temps. Ainsi $E = 1$ signifie \emph{la production a monté d'exactement
ce que les entités traitées produisent de plus, ni plus ni moins} ;
$E = 2{,}5$ signifie qu'elle a monté deux fois et demie plus que cela, donc
que les entités \emph{non} traitées en ont profité aussi --- c'est ce qu'on
appelle ici un \textbf{rebond}. $E < 1$ signale au contraire que le système
a absorbé une partie du gain.

\paragraph{Deux pièges de dénominateur, tous deux corrigés.} D'abord,
\code{tech\_series} mesure la part de production \emph{après} application
du levier ; si $p$ est la part ex ante, la part observée vaut
$s = mp/(1+(m-1)p)$. On inverse cette relation ($p = s/(m-(m-1)s)$) avant
tout calcul : sur le bras ci-dessus, $s = 27{,}08\,\%$ alors que
$p = 19{,}84\,\%$ ; utiliser $s$ directement gonflerait le dénominateur de
$36\,\%$ à l'horizon 1, et diviserait donc $E$ par $1{,}36$. Ensuite, sur un levier en
$\gamma$ l'amplitude n'est pas le rapport des exposants mais
$K^{\Delta\gamma}$, qui dépend du capital ; elle est donc \emph{mesurée} et
non imposée. \fait{} Ces deux corrections ne reposent plus sur une
inversion : $m$ et $p$ sont désormais mesurés \emph{entité par entité} au
pas de l'intervention (\autoref{sec:validation}), et la formule d'inversion
est vérifiée par cette mesure au lieu d'être supposée.

\section{Protocole}

\subsection{Fenêtre, horizon, graines}

\fait{} $t_0 = 2000$, fenêtre $W = 2000$ pas, 5 graines par cellule.
Justification lue dans des résultats \emph{déjà sur disque}, sans
remesurer la relaxation (le prompt §7 l'interdit) : les trois graines
baseline de M4.3 rapportent $t_{\text{converge}}$ = 1016 / 948 / 841 pour
\code{int\_in}, leur variable la plus lente
(\code{simulation\_lab\_data/runs/m4\_3\_\_d1\_\_baseline\_\_seed*/analysis.json}),
et \code{prod\_tot} lui-même est stationnaire par blocs de 500 dès
$t \approx 500$ (moyennes de blocs entre 31\,300 et 32\,100 sur les trois
graines). $t_0=2000$ laisse donc un facteur 2 sur la variable la plus
lente. La fenêtre $W=2000$ vaut environ $38\times$ le temps
d'autocorrélation intégré de \code{prod\_tot} (52 pas, mesuré sur la queue
de ces mêmes séries).

\begin{figure}[H]
\centering
\includegraphics[width=\textwidth]{figures/window_choice.png}
\caption{D'où viennent $t_0$ et $W$. À gauche, \code{prod\_tot} moyenné par
blocs de 500 pas sur les trois graines baseline de M4.3 : le niveau est
stationnaire dès $t \approx 500$, bien avant $t_0$, et les traits pointillés
situent le $t_{\text{converge}}$ de la variable la plus lente du modèle. À
droite, l'autocorrélation qui dimensionne la fenêtre. Toutes ces séries
étaient déjà sur disque ; aucune relaxation n'a été remesurée. Données :
\code{results/analysis/window\_choice.json}.}
\label{fig:fenetre}
\end{figure}

\subsection{Neuf bras, un amorçage partagé}

Un seul run $0 \to t_0$ par graine ; tous les bras sont branchés sur son
\emph{snapshot}. Les bras d'une même graine sont donc rigoureusement
identiques jusqu'à $t_0$ et ne divergent qu'à l'intervention.

\begin{center}\footnotesize
\begin{tabular}{@{}llll@{}}
\toprule
\textbf{Bras} & \textbf{Portée} & \textbf{Levier} & \textbf{Rôle} \\
\midrule
\code{control} & --- & --- & référence inerte \\
\code{null} & fraction $\varphi=0{,}2$ & aucun & plancher de bruit \\
\code{frac\_A150\_phi20} & fraction $\varphi=0{,}2$ & $A\times1{,}5$ & cellule principale \\
\code{frac\_A125\_phi20} & fraction $\varphi=0{,}2$ & $A\times1{,}25$ & échelle d'amplitude \\
\code{frac\_A150\_phi05} & fraction $\varphi=0{,}05$ & $A\times1{,}5$ & échelle d'intensité \\
\code{frac\_A150\_phi50} & fraction $\varphi=0{,}5$ & $A\times1{,}5$ & échelle d'intensité \\
\code{all\_A150} & toutes & $A\times1{,}5$ & hausse globale \\
\code{new\_A150} & nouvelles & $A\times1{,}5$ & vintage technologique \\
\code{frac\_g060\_phi20} & fraction $\varphi=0{,}2$ & $\gamma: 0{,}5\to0{,}6$ & régime (c) du noyau \\
\bottomrule
\end{tabular}
\end{center}

\fait{} 45 runs, tous terminés avec le statut \code{ok}, tous à
$t_{\text{final}} = 4000$, aucun carnet incohérent, aucune transaction
plafonnée (vérifications automatiques dans \code{scripts/analyse.py}, qui
refuse d'analyser une campagne incomplète plutôt que de moyenner des
séries tronquées).

\subsection{L'appariement : pourquoi le bras \code{null} est indispensable}

Le prompt (§2) impose que le tirage de la cohorte \code{fraction} utilise
le générateur de la simulation. Ce tirage \emph{consomme} le flux
aléatoire : dès le pas de l'intervention, un bras \code{fraction} n'a plus
la même suite pseudo-aléatoire qu'un contrôle inerte, et ses naissances
diffèrent immédiatement. Comparer un bras \code{fraction} au contrôle
mélangerait donc le traitement et un simple décalage de flux.

Le bras \code{null} lève exactement cette confusion : il tire la même
cohorte, dans le même état du générateur, et ne lui applique rien (la
valeur assignée est la valeur courante). C'est la référence appariée des
bras \code{fraction}. Les bras \code{all} et \code{new}, qui ne consomment
aucun tirage, sont appariés au contrôle. Le contraste
contrôle\,$\leftrightarrow$\,\code{null} ne contient aucun traitement :
c'est le \textbf{plancher de bruit}.

\fait{} Moyennes de fenêtre du contraste contrôle\,$\leftrightarrow$\,%
\code{null} : $-1{,}01\pm1{,}25$\,\% ($h\le50$), $-0{,}92\pm1{,}08$\,\%,
$-0{,}17\pm0{,}47$\,\%, $+0{,}37\pm0{,}45$\,\% --- toutes à moins de
$1\sigma$ de zéro.\footnote{\label{note:sigma}\textbf{Ce que « $x\,\sigma$ »
veut dire dans ce rapport, et ce que cela ne veut pas dire.} Chaque cellule
est mesurée sur $n = 5$ graines appariées. Pour une quantité $X$ (typiquement
l'écart relatif de \code{prod\_tot} entre un bras et sa référence, moyenné
sur une fenêtre), on dispose des $n$ valeurs par graine $X_1,\dots,X_n$. On
rapporte la moyenne $\bar X$ et l'\textbf{erreur-type}
$\mathrm{SE} = s/\sqrt{n}$ où $s^2 = \frac{1}{n-1}\sum_i (X_i-\bar X)^2$ est
la variance d'échantillon. Ce qu'on note « $x\,\sigma$ » est le rapport
$x = \bar X / \mathrm{SE}$. Comme $\sigma$ est \emph{estimé} sur les mêmes
5 valeurs et non connu, ce rapport ne suit pas une loi normale mais une loi
de \textbf{Student à $n-1 = 4$ degrés de liberté}. Les seuils sont donc
sensiblement plus exigeants que les seuils normaux familiers : pour rejeter
« l'effet est nul » au risque 5\,\% il faut $|x| \ge 2{,}78$ et non
$1{,}96$ ; au risque 1\,\%, $|x| \ge 4{,}60$ et non $2{,}58$. Une mesure à
$3{,}2\,\sigma$ est significative à 5\,\% mais \emph{pas} à 1\,\%
($p \approx 3{,}3\,\%$), là où la lecture normale suggérerait $0{,}1\,\%$.
Une mesure à $5{,}4\,\sigma$ l'est à 1\,\% ($p \approx 0{,}6\,\%$). Ces
seuils s'appliquent à chaque chiffre en $\sigma$ du rapport. L'appariement
par graine est ce qui rend $n=5$ utilisable : il élimine la variance
commune aux deux bras, et l'incertitude portée ici est celle du
\emph{contraste}, pas celle du niveau.} Le décalage de flux n'introduit donc aucun biais
systématique, mais il fixe une résolution : une moyenne de fenêtre plus
petite que $\sim$1\,\% de \code{prod\_tot} n'est pas interprétable
(\autoref{fig:noise}). Sur un pas isolé, la bande est de
\input{tables/noise_floor.tex}\unskip.

\begin{figure}[H]
\centering
\includegraphics[width=\textwidth]{figures/noise_floor.png}
\caption{Plancher de bruit : écart du bras \code{null} au contrôle, cinq
graines. Aucun traitement n'est appliqué ; tout l'écart vient du décalage
de flux aléatoire dû au tirage de la cohorte.}
\label{fig:noise}
\end{figure}

\subsection{Validation de la chaîne de mesure}
\label{sec:validation}

\paragraph{Pourquoi l'horizon 1 est un test exact, et non approché.}
À l'horizon 1, l'écart relatif observé doit valoir \emph{exactement} l'effet
mécanique $(m-1)p$. L'argument est celui-ci, et il mérite d'être déroulé
parce qu'il n'est pas évident : le choc multiplicatif frappe \emph{après}
l'intervention et \emph{avant} la production, ce qui pourrait laisser croire
que l'égalité n'est vraie qu'en moyenne, à la loi des grands nombres près.
Elle est en fait exacte, pour une raison d'appariement. L'intervention ne
touche que $(A,\gamma)$ : au pas $h=1$, l'ensemble des vivantes, son ordre
d'itération et l'état du générateur sont \emph{identiques} dans le bras
traité et dans sa référence appariée. Le même vecteur de chocs $\xi$ frappe
donc les mêmes capitaux dans les deux bras, et les capitaux post-choc
$K_i e^{\xi_i}$ sont eux aussi identiques. Le rapport
$\sum A' (K_ie^{\xi_i})^{\gamma'} / \sum A (K_ie^{\xi_i})^{\gamma}$ est alors
l'amplitude exacte de ce pas-là : le choc ne se simplifie pas terme à terme,
il est simplement \emph{commun aux deux bras}, ce qui suffit. C'est un test,
pas un résultat.

\paragraph{$m$ et $p$ sont désormais mesurés, pas reconstruits.} \fait{} La
première version de ce rapport déduisait l'amplitude et la part traitée des
seules séries agrégées, en inversant $s = mp/(1+(m-1)p)$ --- une
reconstruction correcte, mais qui suppose la formule qu'elle sert à
vérifier. Les deux quantités sont maintenant mesurées \emph{entité par
entité} (\code{scripts/exact\_amplitude.py}). Le moteur n'a pas été modifié
pour cela : après le pas $t_0{+}1$, la production individuelle
$\mathrm{prod}_i = A' K_i^{\gamma'}$ est en mémoire, donc le capital au
moment de produire se retrouve exactement par
$K_i = (\mathrm{prod}_i/A')^{1/\gamma'}$ --- un aller-retour flottant, pas
une approximation --- et l'on peut écrire la production
\emph{contrefactuelle} $A K_i^{\gamma}$ de chaque entité traitée sous son
ancienne technologie.

\input{tables/exact_amplitude.tex}

\begin{figure}[H]
\centering
\includegraphics[width=\textwidth]{figures/calibration_h1.png}
\caption{Calibration de la chaîne de mesure. À gauche, l'écart relatif
observé à l'horizon 1 contre l'effet mécanique prédit : les sept bras
s'alignent sur la première bissectrice sur deux ordres de grandeur d'effet.
À droite, l'amplitude du levier mesurée entité par entité, comparée à celle
imposée. Le bras $\gamma$ n'a pas d'amplitude imposée : la sienne vaut
$K^{\Delta\gamma}$, mesurée à $1{,}9462$.
Données : \code{results/analysis/exact\_amplitude.csv}.}
\label{fig:calibration}
\end{figure}

\fait{} Les amplitudes mesurées valent $1{,}500000$, $1{,}250000$,
$1{,}500000$, $1{,}500000$ et $1{,}500000$ pour les cinq bras à levier sur
$A$ --- les valeurs imposées, à six décimales --- et $1{,}000000$ pour le
bras \code{null}, qui n'applique rien. \fait{} Le bras \code{all\_A150}
donne une part traitée de $1{,}000000$ exactement : les entités nées au pas
de l'intervention reçoivent elles aussi la nouvelle technologie par défaut,
et aucune production n'échappe au traitement. \fait{} Le bras $\gamma$ donne
une amplitude effective de $1{,}946184$ : le passage
$\gamma: 0{,}5 \to 0{,}6$ n'est \emph{pas} un choc de même taille que
$A\times1{,}5$, il est presque deux fois plus fort, ce qui interdit de le
comparer aux bras $A$ comme s'il s'agissait de la même amplitude.

\fait{} \textbf{L'ancienne reconstruction est validée par la mesure} : elle
donnait $1{,}500001$, $1{,}250000$, $1{,}500002$, $1{,}500001$ et
$1{,}946203$ --- à $1{,}5\times10^{-5}$ près en relatif de la mesure directe.
L'inversion ex post / ex ante était donc juste ; elle n'est simplement plus
nécessaire. \fait{} Et le test lui-même passe à la précision machine :
$E(h{=}1) = 1$ sur les six bras à levier, avec un écart maximal à 1 de
$1{,}0\times10^{-15}$. L'exactitude annoncée au paragraphe précédent est
donc vérifiée, pas seulement argumentée.

\section{Résultats}

\input{tables/campaign_windows.tex}

\subsection{Portée partielle, court horizon : rebond observé}

\fait{} Sur la fenêtre d'impact ($h\le50$), tous les bras \code{fraction}
répondent nettement plus fort que la proportionnalité :

\begin{center}
\begin{tabular}{@{}lrrr@{}}
\toprule
\textbf{Bras} & \textbf{écart} & \textbf{significativité} & \textbf{élasticité $E$} \\
\midrule
\code{frac\_A150\_phi05} & $+6{,}84 \pm 2{,}17$\,\% & $3{,}2\sigma$ & $2{,}40$ \\
\code{frac\_A125\_phi20} & $+13{,}02 \pm 1{,}46$\,\% & $8{,}9\sigma$ & $2{,}86$ \\
\code{frac\_A150\_phi20} & $+24{,}73 \pm 1{,}19$\,\% & $20{,}7\sigma$ & $2{,}32$ \\
\code{frac\_A150\_phi50} & $+65{,}46 \pm 2{,}65$\,\% & $24{,}7\sigma$ & $2{,}78$ \\
\code{frac\_g060\_phi20} & $+66{,}79 \pm 1{,}40$\,\% & $47{,}7\sigma$ & $2{,}44$ \\
\bottomrule
\end{tabular}
\end{center}

\fait{} $E$ reste compris entre $2{,}3$ et $2{,}9$ à travers deux
amplitudes ($\times1{,}25$ et $\times1{,}5$), trois intensités
($\varphi = 0{,}05$, $0{,}2$, $0{,}5$) et deux leviers ($A$ et $\gamma$),
sans tendance monotone ni en $\varphi$ ni en amplitude.
\incertitude{} La dispersion de $E$ dans cette plage n'est pas
interprétable en l'état : la cellule $\varphi=0{,}05$, la plus basse en
intensité, n'est mesurée qu'à $3{,}2\sigma$, ce qui, lu correctement en
Student à 4 degrés de liberté (note~\ref{note:sigma}), vaut $p = 3{,}3\,\%$
--- significatif au risque 5\,\%, mais pas à 1\,\%. Ce qui est solide est le
\emph{signe et l'ordre de grandeur} --- une amplification d'un facteur
$\approx$2,5 --- pas son classement d'une cellule à l'autre. Le pic n'est pas au bord de la
fenêtre --- la \autoref{fig:response} montre qu'il est atteint vers
$h \approx 30$ pour les bras $A$ et vers $h \approx 100$ pour le bras
$\gamma$, ce dernier culminant à $+130$\,\%.

\begin{figure}[H]
\centering
\includegraphics[width=\textwidth]{figures/response_prod.png}
\caption{Réponse appariée de \code{prod\_tot}. En haut, portée
\code{fraction} contre le bras \code{null} ; en bas, portées globales
contre le contrôle. La bande grise est le plancher de bruit. Les bras
partiels reviennent dans la bande vers $h\approx250$ ; les bras globaux
s'installent sur un plateau très au-dessus.}
\label{fig:response}
\end{figure}

\subsection{Portée partielle, long horizon : non tranchable}

\fait{} Au-delà de $h\approx300$, tous les bras \code{fraction} à levier
sur $A$ rentrent dans la bande de bruit : $-1{,}8\sigma$ à $+0{,}7\sigma$
sur les fenêtres « moyen » et « établi ». \textbf{Ce ne sont pas des
effets négatifs : ce sont des non-résultats.}

La raison est structurelle et se lit directement sur l'intensité de
traitement observée (\autoref{fig:share}, panneau de gauche) : la part de
production détenue par la cohorte traitée passe de $0{,}198$ (ex ante, à
$h{=}1$) à $0{,}166$, puis $0{,}055$, puis $0{,}022$ sur les quatre
fenêtres. La cohorte s'éteint par renouvellement de la population, et
\emph{aucune naissance ne la reconstitue} --- c'est la sémantique même de
la portée \code{fraction}, qui ne modifie pas la valeur par défaut à la
naissance. Le dénominateur de \eqref{eq:elasticite} tend vers zéro ;
l'élasticité y devient un rapport $0/0$, ce qui est la raison pour
laquelle les courbes de la \autoref{fig:share} sont coupées dès que la
réponse proportionnelle attendue passe sous 2\,\% de \code{prod\_tot}.

\paragraph{Le retour à zéro passe par un creux, et ce creux se lit.}
Rappel de ce qu'on regarde : pour chaque bras \code{fraction}, la courbe est
l'écart relatif de \code{prod\_tot} entre le bras et sa référence appariée,
tracé en fonction de l'horizon $h$. Elle part de l'effet mécanique à $h=1$,
monte jusqu'à un pic vers $h\approx30$, puis redescend à mesure que la
cohorte traitée s'éteint. La question est : \emph{comment} redescend-elle ?

\fait{} Elle ne redescend pas jusqu'à zéro pour s'y arrêter --- elle
\textbf{passe en dessous}. Les deux bras les plus intenses atteignent un
minimum franchement négatif avant de remonter : $-7{,}76 \pm 2{,}01$\,\% à
$h = 195$ pour $\varphi = 0{,}5$, et $-12{,}98 \pm 2{,}40$\,\% à $h = 303$
pour le bras $\gamma$ (soit $3{,}9\sigma$ et $5{,}4\sigma$,
note~\ref{note:sigma} : significatifs à 5\,\% tous les deux, à 1\,\%
seulement pour le second). Autrement dit, pendant plusieurs centaines de
pas, un système qu'on a rendu \emph{plus} productif produit \emph{moins}
qu'un système qu'on n'a pas touché.

\fait{} Ce que montrent les autres agrégats au fond du creux :

\begin{center}\footnotesize
\begin{tabular}{@{}lrr@{}}
\toprule
au fond du creux & $\varphi = 0{,}5$ ($h=195$) & bras $\gamma$ ($h=303$) \\
\midrule
production \code{prod\_tot} & $-7{,}8$\,\% & $-13{,}0$\,\% \\
population \code{pop}       & $-20{,}7$\,\% & $-39{,}1$\,\% \\
capital agrégé \code{K\_tot} & $-9{,}1$\,\% & $-5{,}9$\,\% \\
\bottomrule
\end{tabular}
\end{center}

\inference{} La population s'effondre bien plus que le capital : il y a donc
beaucoup moins d'entités, mais chacune est bien plus riche. Ce n'est pas une
pénurie de capital, c'est une pénurie d'\emph{entités}. L'explication
disponible est d'échelle : pendant le pic, le capital moyen par entité a
monté, alors que le capital de naissance $K_0$ est resté à sa valeur de
départ ; les entrantes naissent donc relativement plus pauvres dans un
système devenu plus grand, et une partie d'entre elles ne survit pas assez
longtemps pour compenser. Le §\ref{sec:ablation} isole ce mécanisme en
régime établi et le teste directement --- ici il n'est qu'une lecture
transitoire, cohérente avec ce qui suit.

\paragraph{Une prémisse du prompt que les données contredisent.} Le §2 du
prompt retient la portée \code{fraction} pour le protocole principal au
motif que \code{new} « a une intensité de traitement qui dérive dans le
temps » alors que \code{fraction} « fixe l'intensité au moment de la
mesure ». \fait{} Dans ce modèle, les deux dérivent : \code{new} monte de
$0$ à $1$ en $\sim$600 pas, \code{fraction} descend de $\varphi$ vers $0$
en $\sim$300 pas. \inference{} La portée \code{fraction} ne répond donc à
la question que sur les premières centaines de pas ; c'est \code{all} et
\code{new}, dont l'intensité converge vers 1, qui portent la version à
intensité fixe de la question. Le protocole a été conservé tel quel --- il
donne les deux réponses --- mais la lecture des résultats doit tenir compte
de cette correction.

\begin{figure}[H]
\centering
\includegraphics[width=\textwidth]{figures/share_elasticity.png}
\caption{À gauche, intensité de traitement réellement observée (part de
\code{prod\_tot} détenue par les entités traitées). À droite, élasticité
normalisée ; la ligne pointillée est la réponse strictement
proportionnelle.}
\label{fig:share}
\end{figure}

\fait{} Une exception : le bras $\gamma$ reste à $+2{,}30 \pm 0{,}30$\,\%
($7{,}6\sigma$) en régime établi, avec une cohorte encore à $0{,}029$ de
part ex ante. \inference{} Les entités à $\gamma$ élevé sont d'autant plus
avantagées que leur capital est grand ($K^{0{,}1}$ croissant) ; leur
cohorte se maintient mieux, et son amplitude effective \emph{augmente} avec
le temps. \incertitude{} C'est aussi ce qui rend l'élasticité de ce bras à
long horizon ($0{,}83$) la moins fiable du tableau : son dénominateur
suppose une amplitude gelée à sa valeur de $h{=}1$.

\subsection{Portée globale, régime établi : réponse sous-proportionnelle}

C'est la cellule où l'intensité de traitement est fixe et vaut 1, donc
celle où l'élasticité est la mieux définie.

\fait{} Une hausse uniforme de 50\,\% de $A$ donne, en régime établi
($h\in[1001,2000]$) :
$+36{,}06 \pm 0{,}41$\,\% ($87{,}9\sigma$) pour la portée \code{all}, et
$+36{,}53 \pm 0{,}53$\,\% ($69{,}3\sigma$) pour la portée \code{new}. Les
deux portées convergent, comme attendu, puisque leur population finit
entièrement traitée ; elles ne diffèrent que par le transitoire. Soit une
élasticité normalisée de $0{,}72$--$0{,}73$, et en élasticité
logarithmique :

\input{tables/campaign_global.tex}

\paragraph{D'où vient ce $0{,}72$ : le calcul, déroulé.} La production
agrégée est le produit de deux facteurs, l'effectif et la production
moyenne : $\mathrm{prod\_tot} = \mathrm{pop} \times
\overline{\mathrm{prod}}$. Le levier agit sur les deux, en sens
\emph{contraires}. \fait{} Mesuré, bras rapporté à son contrôle apparié :

\begin{center}\footnotesize
\begin{tabular}{@{}lrl@{}}
\toprule
grandeur & rapport bras/contrôle & \\
\midrule
population \code{pop} & $\times\,0{,}754$ & il y a moins d'entités \\
production par entité $\overline{\mathrm{prod}}$ & $\times\,1{,}802$
  & chacune produit plus \\
\midrule
production agrégée \code{prod\_tot} & $\times\,1{,}359$
  & $= 0{,}754 \times 1{,}802$ \\
\bottomrule
\end{tabular}
\end{center}

\noindent (rapports mesurés sur les bras d'ablation, qui portent les mêmes
interventions que \code{all\_A150} avec les diagnostics de mortalité en
plus ; fenêtre $h \in [1001,2000]$.) Le $\times1{,}359$ agrégé est donc
exactement le produit des deux mouvements, et la sous-proportionnalité
($1{,}359 < 1{,}500$) vient entièrement de la contraction de population.

\inference{} Reste à comprendre pourquoi une entité produit $\times1{,}802$
alors qu'on n'a multiplié $A$ que par $1{,}5$. La réponse est qu'elle n'a
pas gardé le même capital. Sa production vaut $A K^{\gamma}$ ; le levier
multiplie $A$ par $1{,}50$, et le capital moyen par entité a lui aussi monté.
De combien ? Le capital agrégé fait $\times1{,}090$ et la population
$\times0{,}754$, donc le capital \emph{par entité} fait
$1{,}090/0{,}754 = 1{,}446$ (773 joules par entité dans le contrôle, 1118
dans le bras). La production par entité doit alors faire
\[
\underbrace{1{,}50}_{\text{effet de } A}
\;\times\;
\underbrace{1{,}446^{\gamma}}_{\text{effet du capital, } \gamma = 0{,}5}
= 1{,}50 \times \sqrt{1{,}446} = 1{,}50 \times 1{,}202 = 1{,}804 ,
\]
à comparer au $1{,}802$ mesuré : le compte y est à $0{,}1$\,\% près. La
sur-réponse individuelle n'a donc rien de mystérieux --- c'est le levier,
plus l'accumulation de capital qu'il déclenche.

\begin{figure}[H]
\centering
\includegraphics[width=\textwidth]{figures/levels_established.png}
\caption{Régime établi : écart au bras apparié pour chaque bras.}
\label{fig:levels}
\end{figure}

\subsection{Canaux observés}

Le prompt (§1) donne trois canaux comme grille de lecture, sans en faire
un livrable. Les mesures suivantes les documentent ; elles ne constituent
pas une attribution causale.

\paragraph{Production directe et accumulation de capital.} La production
s'ajoute au capital de l'entité \emph{avant} la phase de marché du même
pas. \fait{} Sur \code{all\_A150} en régime établi, le capital par entité
monte de 45\,\% --- soit une avancée partielle vers la nouvelle échelle
autarcique, qui passe de $\Kaut = 9801$ à $22\,052$ quand $A$ passe de
$1{,}0$ à $1{,}5$ ($\Kaut = [A(1-\delta)/\delta]^{1/(1-\gamma)}$).

\paragraph{Échanges : le sens du prêt éteint des paires entières.}
\fait{} Dans les bras à population hétérogène, le compteur
\code{mkt\_blocked\_dir} --- paires où l'optimum de production jointe
voudrait faire circuler le capital du pauvre vers le riche, ce que le
marché interdit --- atteint 22,6\,\% des rounds (\code{frac\_A150\_phi20},
$h\le50$), 37,3\,\% ($\varphi=0{,}5$) et 36,4\,\% (\code{new\_A150},
$h\le50$), et redescend vers zéro à mesure que la population
s'homogénéise. \fait{} Sur \code{all\_A150}, population homogène, ce
compteur est \emph{exactement} nul --- somme rigoureusement nulle sur les
$5 \times 2000 = 10\,000$ lignes de série du bras. \inference{} Ce canal
est donc une conséquence de l'\emph{hétérogénéité}, pas du niveau de $A$ :
les entités dopées accumulent, deviennent le côté riche de leurs paires, et
cessent alors de prêter.

\paragraph{Faillites.} Les taux ci-dessous sont donnés dans les deux bases,
parce que la population du bras traité est 25\,\% plus petite : lus en
valeur absolue, ils paraissent stables ; rapportés à la population, ils ne
le sont pas du tout.

\fait{} Sur \code{all\_A150} en régime établi :

\begin{center}
\begin{tabular}{@{}lrrr@{}}
\toprule
& \textbf{contrôle} & \textbf{\code{all\_A150}} & \textbf{rapport} \\
\midrule
morts par pas (absolu) & $29{,}96$ & $29{,}99$ & $\times1{,}001$ \\
morts par pas et par 1000 entités & $26{,}50$ & $35{,}19$ & $\times1{,}328$ \\
racines d'insolvabilité par pas (absolu) & $6{,}37$ & $7{,}36$ & $\times1{,}155$ \\
racines d'insolvabilité par 1000 entités & $5{,}64$ & $8{,}64$ & $\times1{,}532$ \\
\bottomrule
\end{tabular}
\end{center}

\inference{} Le nombre absolu de morts par pas est fixé par le flux de
naissances : à l'état stationnaire il vaut $\lambda=30$ quoi qu'il arrive.
Ce qui change réellement est le \emph{taux} : un tiers de morts en plus par
entité et par pas, et une insolvabilité en hausse de 53\,\% par entité --- donc
des durées de vie plus courtes, ce qui est exactement l'autre face de la
contraction de population.

\paragraph{Une explication séduisante, et fausse.} Le candidat immédiat est
un service d'intérêts alourdi : le taux dérive du rendement marginal
$A\gamma K^{\gamma-1}$, proportionnel à $A$, et le service rapporté au
capital monte effectivement de $0{,}0407$ à $0{,}0491$ ($+21$\,\%). Mais
c'est le mauvais dénominateur. \fait{} Rapporté à ce que les entités
\emph{produisent} --- la grandeur qui décide si elles peuvent payer --- le
service \textbf{baisse} : $1{,}128$ dans le contrôle contre $1{,}093$ dans le
bras traité. Le service est plus léger, pas plus lourd. \inference{}
L'attribution au service d'intérêts est donc réfutée par les données mêmes
qui semblaient la soutenir ; elle figurait dans une version antérieure de ce
rapport, elle est retirée. Le §\ref{sec:ablation} traite la question par une
ablation dédiée.

\begin{figure}[H]
\centering
\includegraphics[width=\textwidth]{figures/service_denominator.png}
\caption{Pourquoi l'explication par le service d'intérêts ne tient pas. Le
même service, deux dénominateurs, deux conclusions opposées : rapporté au
capital il monte de 21~\%, rapporté à la production --- ce que l'entité peut
effectivement payer --- il baisse de 3~\%. Données :
\code{results/analysis/service\_denominator.json}.}
\label{fig:denominateur}
\end{figure}

\fait{} Les autres écarts du bras global, pour mémoire : le nombre de
contrats actifs tombe à $\times0{,}602$ et le nombre de prêts par entité de
$30{,}4$ à $24{,}3$ ; le surplus coopératif créé par le marché fait
$\times1{,}92$.

\begin{figure}[H]
\centering
\includegraphics[width=\textwidth]{figures/channels.png}
\caption{Trois canaux. \textbf{Volume de prêt} : les portées globales
s'installent à $+35$\,\%, les portées partielles reviennent à zéro.
\textbf{Paires refusées par le sens du prêt} : le canal n'existe que sous
hétérogénéité et s'éteint avec elle ; il disparaîtra avec la règle qui le
crée, mais tant qu'elle est là il est réel. \textbf{Capital agrégé} :
l'échelle est coupée à $+0{,}7$ parce que tout se joue en dessous ; l'encart
donne la vue d'ensemble, où l'on voit le seul dépassement, celui du bras
$\gamma$ pendant son transitoire. Le panneau « morts par pas » de la version
précédente a été retiré : à l'état stationnaire le flux de morts est fixé
par le flux de naissances et vaut $\lambda$ quoi qu'il arrive --- c'est le
\emph{taux} qui porte l'information, et il est traité au §\ref{sec:ablation}
et à la \autoref{fig:tension_mortalite}.}
\label{fig:channels}
\end{figure}

\subsection{Une statistique qui ne dépend d'aucune fenêtre}

\paragraph{Le problème.} Toutes les élasticités données jusqu'ici sont des
moyennes sur une fenêtre d'horizons ($h\le50$, $h\in[1001,2000]$\dots). Le
choix de ces fenêtres est défendable, mais il reste un choix : un lecteur
peut légitimement demander ce que devient le résultat si l'on découpe
autrement. On donne donc une seconde statistique qui n'a pas de fenêtre du
tout.

\paragraph{L'idée, en mots.} Au lieu de comparer une moyenne de réponses à
une moyenne de prédictions, on additionne \emph{tous} les joules gagnés sur
tout l'horizon, on additionne \emph{tous} les joules que la
proportionnalité stricte aurait prédits sur le même horizon, et on divise
l'un par l'autre :
\begin{equation}
E_{\text{cum}} \;=\;
\frac{\text{total des joules effectivement gagnés}}
     {\text{total des joules prédits par la proportionnalité}}
\;=\;
\frac{\displaystyle\sum_{h} \Delta\mathrm{prod}(h)}
     {\displaystyle\sum_{h} (m-1)\,p(h)\,\mathrm{prod}_{\text{réf}}(h)} .
\label{eq:cumule}
\end{equation}
La somme au numérateur est l'aire sous la courbe de réponse ; celle au
dénominateur est l'aire sous la courbe de la réponse proportionnelle. Leur
rapport se lit comme les $E$ de fenêtre : $1$ = proportionnel, $>1$ =
rebond, $<1$ = absorbé. Aucun découpage n'y intervient, et les horizons où
l'effet est gros pèsent naturellement plus que ceux où il est petit.

\fait{} Sur les bras à intensité de traitement stable :
$E_{\text{cum}} = 0{,}744 \pm 0{,}002$ (\code{all\_A150}) et
$0{,}765 \pm 0{,}008$ (\code{new\_A150}) --- soit environ trois quarts de la
réponse proportionnelle, en accord avec les $0{,}72$--$0{,}73$ des fenêtres
établies. \fait{} Sur les bras partiels : $1{,}116 \pm 0{,}050$
(\code{frac\_g060\_phi20}) et $0{,}638 \pm 0{,}082$
(\code{frac\_A150\_phi50}).

\incertitude{} \textbf{Cette statistique cesse d'être exploitable dès que la
cohorte traitée s'éteint.} Pour les bras \code{fraction} de faible
intensité, la longue traîne d'horizons où $p(h) \approx 0$ contribue
presque zéro au dénominateur mais du bruit au numérateur :
$E_{\text{cum}} = 0{,}04 \pm 0{,}13$ pour $\varphi=0{,}2$, un chiffre qui ne
mesure rien. \inference{} $E_{\text{cum}}$ confirme donc les élasticités de
fenêtre là où celles-ci sont définies --- portées globales, et bras partiels
tant que la cohorte pèse --- sans les remplacer, et sans rien sauver là où
elles ne le sont pas.

\section{Ablation : d'où vient la contraction de population ?}
\label{sec:ablation}

Le §3.3 laisse une question ouverte, et le §3.5 vient d'en écarter la
réponse la plus naturelle. Il reste une hypothèse, désignée par les données
elles-mêmes : le capital de naissance $K_0$ est \emph{fixe} alors que
l'échelle du système bouge. \fait{} Une entité neuve vaut 3,23~\% du capital
moyen dans le contrôle, contre 2,24~\% dans le bras traité --- elle naît
relativement plus pauvre, dans un système plus grand.

\subsection{Plan de l'ablation}

Cinq bras, mêmes graines et mêmes snapshots à $t_0$ que la campagne
principale, donc directement comparables à elle. Le facteur de compensation
$2{,}25 = A^{1/(1-\gamma)} = 1{,}5^{2}$ est celui qui laisse invariant le
rapport $K_0/\Kaut$ ; c'est la convention déjà employée par M4.3 pour ses
cellules compensées. Le facteur $1{,}45$ est le rapport \emph{observé} des
capitaux par entité --- une compensation a posteriori plutôt qu'a priori.

Ces bras enregistrent en plus l'âge au décès de chaque entité, ce que la
campagne principale ne faisait pas, afin de répondre directement à « qui
meurt en plus ? ». \fait{} Le drapeau d'enregistrement n'écrit que dans une
liste : le bras \code{abl\_control} reproduit \code{control} \textbf{bit à
bit} sur six colonnes et 5 graines (écart maximal exactement nul), ce qui
vérifie que la mesure supplémentaire ne perturbe rien.

\subsection{Résultat}

\input{tables/ablation_k0.tex}

\fait{} Compenser $K_0$ à l'échelle autarcique \textbf{annule complètement}
la contraction : la population revient à $\times1{,}005 \pm 0{,}005$ de son
contrôle, et tous les diagnostics de mortalité avec elle (durée de vie
$37{,}9$ contre $37{,}7$ pas, racines d'insolvabilité $\times1{,}018$ contre
$\times1{,}155$ sans compensation ; tableau ci-dessus). Le bras à
compensation partielle ($\times1{,}45$) se place entre les deux sur chacun
de ces indicateurs.

\inference{} Ce retour au contrôle n'est pas une coïncidence numérique :
c'est la manifestation d'une \textbf{covariance d'échelle} du modèle, déjà
connue de la lignée sous la forme « on peut multiplier $A$ et $K_0$ de façon
cohérente sans changer le fonctionnement du système ». Le §\ref{sec:echelle}
la démontre et la met à l'épreuve. Ce qui est nouveau ici, et ce qui compte
pour le verdict, n'est donc pas que la compensation marche, mais que
\emph{sans elle} un simple choc sur $A$ contracte la population de 25\,\% :
$K_0$ n'est pas un paramètre neutre qu'on peut laisser à sa valeur pendant
qu'on bouge la technologie.

\begin{figure}[H]
\centering
\includegraphics[width=\textwidth]{figures/ablation_k0_levels.png}
\caption{Ablation sur le capital de naissance. Dans le panneau de gauche, la
courbe du bras compensé à l'échelle autarcique (rouge) revient se superposer
au contrôle (gris), tandis que le bras à $K_0$ fixe (bleu) s'installe 25~\%
en dessous. Données : \code{results/ablation\_k0/}.}
\label{fig:ablation}
\end{figure}

\begin{figure}[H]
\centering
\includegraphics[width=\textwidth]{figures/ablation_k0_survival.png}
\caption{Âge au décès, fenêtre $h \in [1001, 2000]$, cinq graines cumulées.
Le surcroît de mortalité du bras à $K_0$ fixe est concentré \emph{aux tout
premiers pas de vie} ; au-delà de vingt pas les distributions se
superposent. Compenser $K_0$ ramène la courbe sur celle du contrôle. Barres
d'incertitude de Poisson.}
\label{fig:survie}
\end{figure}

\subsection{Ce que cela change au verdict}

\fait{} Avec $K_0$ compensé, la production agrégée répond
$\times2{,}267 \pm 0{,}011$ à un choc $\times1{,}5$ sur $A$, soit une
élasticité logarithmique $\varepsilon = 2{,}018$. \inference{} Cette valeur
n'est pas quelconque : $1/(1-\gamma) = 2$ exactement. C'est l'exposant de
l'échelle autarcique $\Kaut \propto A^{1/(1-\gamma)}$, et la production y
vaut $A\,K^{\gamma} \propto A \cdot A^{\gamma/(1-\gamma)} = A^{1/(1-\gamma)}$.
Autrement dit : \textbf{quand le capital de naissance suit l'échelle, le
système entier se remet à l'échelle de sa référence autarcique}, et la
réponse est fortement super-proportionnelle --- élasticité normalisée
$E = 2{,}53$.

Les deux chiffres --- $\varepsilon = 0{,}76$ à $K_0$ fixe,
$\varepsilon = 2{,}02$ à $K_0$ compensé --- ne se contredisent pas : ils
répondent à deux questions différentes, et il faut choisir laquelle on pose.

\begin{center}
\begin{tabularx}{\textwidth}{@{}lX@{}}
\toprule
\textbf{Question posée} & \textbf{Réponse mesurée} \\
\midrule
« La technologie des entités s'améliore de 50~\%, mais les entrants restent
dotés comme avant. » & Réponse \textbf{sous}-proportionnelle,
$\varepsilon = 0{,}76$. Le gain par entité est réel ($\times1{,}80$) mais la
population se contracte d'un quart, parce que les nouveaux nés sont devenus
relativement trop petits pour survivre. \\
\addlinespace
« L'économie entière, dotation des entrants comprise, monte d'un cran
technologique. » & Réponse \textbf{super}-proportionnelle,
$\varepsilon = 2{,}02 = 1/(1-\gamma)$. La population est inchangée, et la
production suit exactement l'échelle autarcique. \\
\bottomrule
\end{tabularx}
\end{center}

Cette coïncidence à 0,9~\% près ne vaudrait rien sur un seul point : elle est
mise à l'épreuve en deux autres $\gamma$ au §\ref{sec:loi}, où la loi prédit
des valeurs franchement différentes.

\fait{} Enfin, $K_0$ est un levier puissant en lui-même : le bras
$K_0 \times 2{,}25$ \emph{seul}, à technologie inchangée, fait
$\times1{,}358$ en population et $\times1{,}694$ en production.

\section{La tension : une seule grandeur pour lire l'état du système}
\label{sec:tension}

Les sections précédentes ont fait apparaître la même chose sous trois noms :
un capital de naissance devenu trop petit, une population qui se contracte,
des durées de vie qui raccourcissent. Cette section propose une grandeur
unique qui les résume, et la met à l'épreuve.

\subsection{Définition}

Une entité \emph{isolée} --- sans marché, sans dette --- a un capital
d'équilibre exact. À chaque pas son capital devient
$(K + A K^{\gamma})(1-\delta)$, dont le point fixe est l'échelle autarcique
déjà rencontrée :
\begin{equation}
\Kaut \;=\; \left[\frac{A\,(1-\delta)}{\delta}\right]^{1/(1-\gamma)} .
\label{eq:kaut}
\end{equation}
C'est une référence purement paramétrique : elle ne dépend ni de la
population, ni du crédit, ni de l'histoire du run.

En face, on veut l'échelle à laquelle le système tourne \emph{réellement}.
On la définit par la production. Le \textbf{capital équivalent} $K_{\text{eq}}$
est le capital qu'auraient toutes les entités d'un groupe technologique si,
à production agrégée et effectif identiques, elles étaient toutes au même
capital :
\begin{equation}
\mathrm{prod} \;=\; n\,A\,K_{\text{eq}}^{\gamma}
\qquad\Longleftrightarrow\qquad
K_{\text{eq}} \;=\; \left(\frac{\mathrm{prod}}{n\,A}\right)^{1/\gamma} .
\label{eq:keq}
\end{equation}
La \textbf{tension} est le rapport des deux :
\begin{equation}
T \;=\; \frac{\Kaut}{K_{\text{eq}}} .
\label{eq:tension}
\end{equation}
$T = 1$ : le système tourne à son échelle autarcique. $T > 1$ : il tourne
\emph{en dessous}, le capital y est maintenu bas par tout ce que l'autarcie
ignore --- le service des intérêts, la mortalité, le renouvellement par des
naissances petites. $T < 1$ : il porte plus de capital que la technologie
seule ne le soutiendrait.

\paragraph{Trois conventions, énoncées parce qu'elles ne sont pas neutres.}
\emph{(i)} $\Kaut$ ignore le service des intérêts --- c'est sa définition,
et c'est ce qui fait de $T$ la mesure de tout ce qui n'est pas la
technologie. \emph{(ii)} $K_{\text{eq}}$ \textbf{n'est pas} le capital
moyen : $K^{\gamma}$ étant concave, la moyenne des productions est
inférieure à la production de la moyenne, donc
$K_{\text{eq}} \le \overline{K}$ par Jensen, avec égalité seulement si tous
les capitaux sont égaux. Le rapport $K_{\text{eq}}/\overline{K}$ est donc
lui-même une mesure d'inégalité interne, et il est rapporté à côté de la
tension. \emph{(iii)} Le calcul est fait par groupe technologique, seul
niveau où $A$ et $\gamma$ sont définis ; l'agrégat d'un run à technologies
multiples est la moyenne des tensions pondérée par la production.

\paragraph{Comment elle est calculée, et sur quelles données.} La tension
est un \emph{dérivé} des séries déjà écrites : \code{tech\_series.csv}
fournit $(\mathrm{prod}, n, A, \gamma)$ par groupe et par pas, et
\code{series.csv} les morts. \fait{} Elle a donc pu être calculée sur les 93
runs déjà terminés, sans en relancer un seul et sans toucher au moteur
(\code{scripts/tension.py}, \code{scripts/tension\_figures.py}). Chaque run
porte désormais \code{tension.csv}, \code{tension\_agg.csv} et
\code{figures/tension.png}, consultables dans \code{simulation\_lab}.
\incertitude{} Un point de méthode, reporté dans la colonne \code{basis} de
chaque fichier : \code{tech\_series} donne l'effectif de \emph{fin} de pas,
alors que la production a été calculée avant les morts du pas. L'effectif
producteur est donc reconstitué par $n + \mathrm{deaths}$ --- exact tant
qu'une seule technologie est vivante, donc pour tous les runs de contrôle,
d'amorçage, d'ablation et de loi d'échelle. Quand plusieurs technologies
coexistent, les morts sont imputées au prorata de l'effectif ; sur les bras
\code{fraction} le groupe traité pèse moins de 2\,\% de la production, et
l'erreur est bornée par les $\approx$30 morts du pas.

\subsection{Ce qu'elle vaut, et ce qu'elle sépare}

\fait{} Dans le régime de référence ($\gamma = 0{,}5$, $A = 1$,
$\delta = 0{,}01$, $K_0 = 25$), $\Kaut = 9801$ et $K_{\text{eq}} \approx 735$ :
la tension vaut $T = 13{,}32 \pm 0{,}02$. \textbf{Le système tourne à un
treizième de son échelle autarcique.} \fait{} L'écart de Jensen y est faible
($K_{\text{eq}}/\overline{K} = 0{,}97$) : les capitaux sont proches les uns
des autres, la tension n'est pas un artefact d'inégalité interne.

\input{tables/tension_arms.tex}

\begin{figure}[H]
\centering
\includegraphics[width=\textwidth]{figures/tension_arms.png}
\caption{Tension en régime établi, campagne et ablation. Le trait pointillé
est le contrôle. À gauche, seuls les bras à portée globale déplacent la
tension --- les bras \code{fraction} restent sur le contrôle, leur cohorte
traitée étant devenue négligeable. À droite, l'ablation : $K_0$ compensé à
l'échelle autarcique (rouge) ramène la tension exactement sur le contrôle,
$K_0$ fixe la porte à 21. Données :
\code{results/analysis/tension\_arms.csv}.}
\label{fig:tension_arms}
\end{figure}

\fait{} La compensation de $K_0$ conserve la tension à $0{,}5$\,\% près :
$13{,}26 \pm 0{,}02$ pour \code{abl\_A150\_K0aut} contre
$13{,}32 \pm 0{,}02$ pour son contrôle, alors que le même choc sans
compensation la porte à $21{,}12 \pm 0{,}04$. \inference{} C'est la lecture
la plus directe du §\ref{sec:ablation} : ce que la compensation préserve,
c'est cette grandeur sans dimension, et c'est pourquoi tous les diagnostics
la suivent.

\subsection{Mortalité et durée de vie contre tension}

\fait{} Rapportées à la tension plutôt qu'aux paramètres, la mortalité et la
durée de vie s'ordonnent. Sur les 100 runs à $\gamma = 0{,}5$ et
$\delta = 0{,}01$ --- couvrant une plage de tension de $3{,}1$ à $45{,}6$,
soit un facteur $15$, obtenue par trois interventions différentes ($A$ seul,
$K_0$ seul, $A$ et $K_0$ ensemble) --- le taux de mortalité suit
\begin{equation}
\frac{\text{morts}}{\text{population}\cdot\text{pas}} \;\propto\; T^{0{,}668},
\qquad R^2_{\log\text{-}\log} = 0{,}9885 ,
\label{eq:mortalite}
\end{equation}
avec un écart médian à l'ajustement de $1{,}4\,\%$. \fait{} Et les trois
leviers s'accordent : leurs biais respectifs valent $-0{,}7\,\%$,
$-4{,}3\,\%$ et $+0{,}4\,\%$. \inference{} À $\delta$ fixé, il est donc
indifférent \emph{par quel chemin} on a atteint un niveau de tension : seul
le niveau compte. \fait{} La durée de vie moyenne, qui est l'inverse du
taux, passe de $114{,}9$ pas à $T = 3{,}1$ à $18{,}3$ pas à $T = 45{,}5$.

\fait{} Le cas le plus net est celui de deux bras qui n'ont rien en commun
que leur tension : diviser $K_0$ par 4 ($K_0 = 6{,}25$) et doubler $A$
($A = 2$) --- deux interventions sans rapport --- amènent le système à
$T = 29{,}46$ et $T = 29{,}61$. Leur mortalité vaut $0{,}04282$ et
$0{,}04304$, leur durée de vie $23{,}35$ et $23{,}23$ pas, leur population
701 et 696. \emph{Le système ne se souvient pas de ce qu'on lui a fait ;
il ne connaît que sa tension.} --- tant que $\delta$ ne bouge pas.

\begin{figure}[H]
\centering
\includegraphics[width=\textwidth]{figures/tension_mortality.png}
\caption{Note [22] du relecteur : mortalité contre tension, \textbf{un
panneau par $\gamma$}, un point par run, axes absolus partagés. L'étoile
pleine est la ligne de base du panneau, les étoiles creuses celles des deux
autres $\gamma$ : les trois familles ne travaillent pas autour du même point.
Chaque encart rejoue le panneau en écart à \emph{sa} base, seule façon de
séparer les bras qui se superposent sur les axes absolus. La durée de vie est
l'inverse exact du taux ($\sum \text{pop}/\sum \text{morts}$ contre
$\sum \text{morts}/\sum \text{pop}$) : elle est portée en second axe, à
droite, et non en second nuage. Aucun ajustement n'est tracé --- voir
ci-dessous. Le bras \code{frac\_g060} porte $\gamma = 0{,}5$ en configuration
mais traite $20\,\%$ des entités à $\gamma = 0{,}6$ ; il a son propre
marqueur. Données : \code{results/analysis/tension\_runs.csv}.}
\label{fig:tension_mortalite}
\end{figure}

\incertitude{} \textbf{Pourquoi aucune droite n'est tracée.} Une version
antérieure de cette figure superposait les trois $\gamma$ dans un seul plan
et y traçait quatre ajustements. Les trois ajustements par $\gamma$ étaient
sans objet : à $\gamma = 0{,}4$ et $\gamma = 0{,}6$, les neuf runs
disponibles couvrent une plage de tension de $0{,}6\,\%$ et $1{,}0\,\%$
respectivement --- ce sont des runs à l'échelle compensée, faits pour
\emph{ne pas} déplacer la tension (§\ref{sec:echelle}). Une droite y ajuste
du bruit de graine.

\fait{} Le quatrième ajustement, « tous $\gamma$ confondus », méritait un
examen plus sérieux : avec $R^2 = 0{,}72$ sur une plage de tension de $2{,}8$
à $88{,}9$, il avait l'air d'une loi. Il n'en est pas une, et on peut le
montrer sans rien ajuster. Les trois lignes de base valent
$(T, \text{morts/pop}) = (4{,}25 ; 0{,}01415)$, $(13{,}32 ; 0{,}02650)$ et
$(88{,}56 ; 0{,}04918)$. La droite qui passe par ces \emph{trois points
seuls} a pour exposant
\[
\frac{\ln(0{,}04918/0{,}01415)}{\ln(88{,}56/4{,}25)} = 0{,}410 ,
\]
contre $0{,}396$ pour l'ajustement sur les $127$ runs. \inference{} Les deux
coïncident à $3{,}5\,\%$ : l'ajustement global ne mesure aucune réponse à la
tension, il mesure le déplacement de la base avec $\gamma$. C'est pourquoi la
figure porte les bases des autres $\gamma$ en marqueurs creux --- la droite
qu'on retirait joignait ces trois étoiles, et rien d'autre.

\incertitude{} \textbf{Reste donc une seule relation défendable, et une
réserve décisive.} La loi de puissance en $T^{0{,}668}$ ($R^2 = 0{,}9885$)
n'est établie qu'à $\gamma = 0{,}5$ \emph{et} $\delta$ fixé ; elle est
ajustée et tracée en §\ref{sec:tension_test} et à la
figure~\ref{fig:tension_sweep}, sur l'échantillon qui balaie effectivement la
tension. Extrapolée à $\gamma = 0{,}6$, elle prédirait à $T = 88{,}6$ un taux
de $0{,}090$ quand le mesuré vaut $0{,}049$. Et la restriction « à $\delta$
fixé » n'est pas cosmétique : c'est l'objet de la sous-section suivante.

\subsection{La tension est-elle un paramètre d'état ? Non : le test qui
échoue}
\label{sec:tension_test}

\paragraph{Le test.} Si la tension résume vraiment l'état du système, alors
n'importe quelle façon de la déplacer doit produire la même mortalité. On a
donc appliqué au même état à $t_0 = 2000$, sur les mêmes snapshots
d'amorçage, \textbf{quatre leviers indépendants} : $K_0$ seul (de 3 à 400,
base 25), $\delta$ seul (de $0{,}005$ à $0{,}04$, base $0{,}01$), $A$ seul
(de $0{,}75$ à $2{,}0$), et $A$ et $K_0$ ensemble. Trente-neuf runs
supplémentaires, trois graines chacun. Le levier $\delta$ est le plus
exigeant, parce qu'il entre \emph{dans la définition même} de $\Kaut$
autant que dans la dynamique.

\input{tables/tension_sweep.tex}

\fait{} \textbf{Le test échoue, et il échoue exactement sur $\delta$.} Sur
les quatre leviers réunis (109 runs), l'ajustement en tension tombe à
$R^2 = 0{,}58$, avec un biais de $+36\,\%$ sur les seuls bras $\delta$ ---
jusqu'à $+102\,\%$ sur un point. La raison se lit directement sur les
niveaux : faire varier $\delta$ d'un facteur 8 déplace la tension d'un
facteur $13{,}4$ ($T = 37{,}3$ à $\delta = 0{,}005$, $T = 2{,}8$ à
$\delta = 0{,}04$) et ne déplace le taux de mortalité que de $14\,\%$
($0{,}0254$ à $0{,}0289$). La loi établie à $\delta$ fixé prédirait
$0{,}0524$ et $0{,}0092$ : elle se trompe d'un facteur 2 dans un sens et 3
dans l'autre.

\fait{} Le contre-exemple symétrique du précédent le montre sans aucune
statistique : $\delta = 0{,}04$ donne $T = 2{,}78$ et $K_0 = 400$ donne
$T = 3{,}12$ --- deux tensions presque égales. Leurs mortalités valent
$0{,}0289$ et $0{,}0087$, un facteur $3{,}3$ ; leurs durées de vie $34{,}6$
et $114{,}9$ pas ; leurs populations 1043 et 3427. \emph{À tension égale,
deux systèmes entièrement différents.}

\inference{} \textbf{La tension n'est donc pas un paramètre d'état du
modèle.} C'est un \emph{diagnostic d'échelle}, valide et précis tant que la
dépréciation ne bouge pas. La raison est visible dans
\eqref{eq:kaut} : $\Kaut$ dépend de $\delta$ à la puissance
$1/(1-\gamma) = 2$, soit bien plus fortement que ne dépend d'elle ce qui
tue réellement les entités.

\paragraph{Ce qui, en revanche, aligne les quatre leviers.} \fait{} La même
mortalité, rapportée à l'\textbf{intensité de rotation du crédit} ---
principal transféré au pas, divisé par le capital agrégé --- donne
\begin{equation}
\frac{\text{morts}}{\text{population}\cdot\text{pas}}
\;\propto\;
\left(\frac{\code{loan\_volume}}{\code{K\_tot}}\right)^{1{,}337},
\qquad R^2_{\log\text{-}\log} = 0{,}9982 ,
\label{eq:rotation}
\end{equation}
avec un écart médian de $0{,}41\,\%$, un écart maximal de $4{,}9\,\%$, et un
biais par levier compris entre $-1{,}7\,\%$ et $+0{,}7\,\%$ --- $\delta$
compris.

\begin{figure}[H]
\centering
\includegraphics[width=\textwidth]{figures/tension_sweep.png}
\caption{Le test des quatre leviers. À gauche, la tension : les leviers
d'échelle ($A$, $K_0$) tombent sur une droite, les points $\delta$
(triangles) s'en écartent franchement --- la tension leur attribue une
variation de mortalité qui n'a pas lieu. À droite, la même mortalité contre
la rotation du crédit : les quatre leviers se confondent sur une seule loi
de puissance. Données :
\code{results/analysis/tension\_runs.csv}.}
\label{fig:tension_sweep}
\end{figure}

\incertitude{} \textbf{Ce résultat ne doit pas être sur-interprété, et voici
pourquoi.} La rotation du crédit n'est pas un paramètre : c'est une variable
d'état \emph{endogène}, mesurée sur le même run que la mortalité.
\eqref{eq:rotation} est donc une régularité descriptive entre deux sorties,
pas une loi de contrôle : elle ne dit pas comment obtenir une mortalité
donnée, elle dit que \emph{si} l'on connaît le débit du marché du crédit
rapporté au stock de capital, on connaît la mortalité à $0{,}4\,\%$ près, quel
que soit le paramètre qu'on a bougé pour l'obtenir. \inference{} Ce que cela
suggère --- sans le démontrer --- est que le canal par lequel les paramètres
tuent est le crédit, et que la tension n'est un bon prédicteur que dans la
mesure où elle suit ce débit. Établir le lien dans l'autre sens, des
paramètres vers la rotation, est le prolongement naturel de ce travail et
n'est pas fait ici.

\begin{figure}[H]
\centering
\includegraphics[width=\textwidth]{figures/tension_trajectoires.png}
\caption{Trajectoires de trois bras contrastés. À gauche, la tension dans le
temps : le bras à $K_0$ compensé reste sur le contrôle, celui à $K_0$ fixe
saute à l'intervention. À droite, le même trajet dans le plan
(tension, mortalité) : le système quitte un point d'équilibre, décrit une
boucle transitoire, et s'installe sur un autre.}
\label{fig:tension_traj}
\end{figure}

\subsection{Ce que $A$ fait à la tension : l'exposant n'est pas 1}
\label{sec:tension_vs_A}

\paragraph{L'observation.} À la lecture de la
figure~\ref{fig:tension_traj}, le relecteur note que multiplier $A$ par
$1{,}5$ semble multiplier la tension par $1{,}5$, et demande si cela vaut en
général. \fait{} C'est presque vrai, et l'écart est systématique : le bras
\code{all\_A150} porte la tension de $13{,}32$ à $21{,}12$, soit
$\times 1{,}586$ et non $\times 1{,}5$ --- $5{,}7\,\%$ de trop, toujours dans
le même sens.

\fait{} Sur les quatre valeurs $A \in \{0{,}75;\,1{,}25;\,1{,}5;\,2{,}0\}$
appliquées au même $t_0$ à $K_0$ fixé, la relation est une loi de puissance
propre --- et son exposant n'est pas $1$ :
\[
T \propto A^{1{,}145}, \qquad R^2_{\log\text{-}\log} = 0{,}9999
\qquad (\gamma = 0{,}5) .
\]

\paragraph{L'exposant n'est pas libre.} \fait{} Il se déduit des deux
définitions. De $\Kaut = [A(1-\delta)/\delta]^{1/(1-\gamma)}$ et
$K_{\text{eq}} = (\text{prod}/(n A))^{1/\gamma}$, en dérivant
$\ln T = \ln \Kaut - \ln K_{\text{eq}}$ par rapport à $\ln A$ :
\begin{equation}
\frac{\mathrm{d}\ln T}{\mathrm{d}\ln A}
\;=\;
\frac{1}{1-\gamma} \;-\; \frac{1}{\gamma}\bigl(\varepsilon_{\text{prod}}
- \eta_{\text{pop}} - 1\bigr) ,
\label{eq:exposant_tension}
\end{equation}
où $\varepsilon_{\text{prod}} = \mathrm{d}\ln(\text{prod})/\mathrm{d}\ln A$
et $\eta_{\text{pop}} = \mathrm{d}\ln n/\mathrm{d}\ln A$ se mesurent sur les
mêmes runs. \incertitude{} Il faut être clair sur ce qu'est
\eqref{eq:exposant_tension} : c'est une \textbf{réécriture des définitions},
pas une hypothèse. Le résidu de $10^{-3}$ entre les deux membres du tableau
ci-dessous mesure la propreté de l'aller-retour numérique
($\text{prod} = n A K_{\text{eq}}^{\gamma}$ inverse exactement la définition
de $K_{\text{eq}}$), et non la validité de l'identité --- qui ne pouvait pas
être fausse.

\paragraph{Ce qui rend le résultat non trivial.} \fait{} Les deux élasticités
sont \emph{presque insensibles à $\gamma$} :
$\varepsilon_{\text{prod}}$ vaut $0{,}64$, $0{,}71$ et $0{,}75$, et
$\eta_{\text{pop}}$ vaut $-0{,}72$, $-0{,}72$ et $-0{,}75$, pour
$\gamma = 0{,}4$, $0{,}5$ et $0{,}6$. \inference{} Toute la dépendance en
$\gamma$ de l'exposant vient donc des deux facteurs \emph{géométriques}
$1/(1-\gamma)$ et $1/\gamma$ de \eqref{eq:exposant_tension}, et non d'un
changement de comportement du système. C'est ce qui la rend prévisible.

\input{tables/tension_vs_A.tex}

\begin{figure}[H]
\centering
\includegraphics[width=\textwidth]{figures/tension_vs_A.png}
\caption{À gauche, la tension contre $A$ à $K_0$ fixé, un trait par $\gamma$,
rapportée à la ligne de base de chaque famille ; le tireté est la droite de
pente 1 que l'observation supposait. Le carré est le contre-point décisif :
la même hausse $A \times 1{,}5$ \emph{avec} $K_0$ compensé ne déplace pas la
tension. À droite, l'exposant mesuré contre celui que donne
\eqref{eq:exposant_tension}. Données :
\code{results/analysis/tension\_vs\_A.csv}.}
\label{fig:tension_vs_A}
\end{figure}

\paragraph{Réponse.} \fait{} Non, $A \times x$ n'implique pas $T \times x$.
L'exposant vaut $0{,}776$, $1{,}145$ et $1{,}673$ pour
$\gamma = 0{,}4$, $0{,}5$ et $0{,}6$ --- il double presque sur cette plage,
avec un $R^2$ d'au moins $0{,}9999$ dans les trois cas. \inference{}
Il ne passe par $1$ qu'au voisinage de $\gamma \approx 0{,}46$ : la
coïncidence observée est une propriété du $\gamma$ de référence, pas du
modèle. À $\gamma = 0{,}6$, $A \times 2$ multiplie la tension par $3{,}21$.

\fait{} \textbf{Et c'est entièrement un phénomène à $K_0$ fixé.} Sous
compensation, la covariance d'échelle (§\ref{sec:echelle}) impose
$\varepsilon_{\text{prod}} = 1/(1-\gamma)$ et $\eta_{\text{pop}} = 0$ ;
\eqref{eq:exposant_tension} donne alors
\[
\frac{1}{1-\gamma} - \frac{1}{\gamma}\left(\frac{1}{1-\gamma} - 1\right)
= \frac{1}{1-\gamma} - \frac{1}{1-\gamma} = 0 .
\]
Mesuré : $-0{,}012$ pour \code{abl\_A150\_K0aut}, contre $+1{,}137$ pour le
même choc sans compensation. \inference{} L'apparent « $A\times x
\Rightarrow T \times x$ » n'est donc pas une propriété de $A$ : c'est une
propriété du \emph{décalage} entre l'échelle du système et un capital de
naissance qu'on a laissé derrière.

\incertitude{} \textbf{Ce que la prédiction a et n'a pas attrapé.} L'exposant
attendu avait été écrit avant de lancer les runs, en supposant
$\varepsilon_{\text{prod}} - \eta_{\text{pop}} \approx 1{,}43$ constant :
$0{,}59$ à $\gamma = 0{,}4$ et $1{,}78$ à $\gamma = 0{,}6$. Le sens est
confirmé --- l'exposant s'éloigne franchement de 1 des deux côtés --- mais
les valeurs sont respectivement $0{,}776$ et $1{,}673$, parce que
$\varepsilon_{\text{prod}} - \eta_{\text{pop}}$ n'est pas tout à fait
constant : il vaut $1{,}356$, $1{,}433$ et $1{,}501$. \inference{} Cette
lente dérive n'est pas expliquée ici, et c'est la seule chose que ce test
laisse ouverte.

\section{La loi d'échelle $\varepsilon = 1/(1-\gamma)$}
\label{sec:loi}

\subsection{Pourquoi cet exposant : la covariance d'échelle}
\label{sec:echelle}

L'exposant $1/(1-\gamma)$ n'est pas un ajustement empirique : il se déduit
du modèle, et l'argument tient en une page. \textbf{Le modèle est covariant
d'échelle.}

\paragraph{L'énoncé.} Posons
\begin{equation}
c \;=\; \left(\frac{A'}{A}\right)^{1/(1-\gamma)} ,
\label{eq:facteur}
\end{equation}
et remplaçons \emph{simultanément} $A \to A'$ et $K_0 \to c\,K_0$, en
laissant $\gamma$, $\delta$, $\sigma$, $\lambda$ et la graine inchangés.
Alors la simulation obtenue est \emph{exactement} la simulation d'origine
avec tous les capitaux multipliés par $c$.

\paragraph{La vérification, phase par phase.} Il suffit de suivre un capital
$K \mapsto cK$ à travers les huit phases d'un pas.

\begin{center}\footnotesize
\begin{tabularx}{\textwidth}{@{}lX@{}}
\toprule
\textbf{Phase} & \textbf{Effet du changement d'échelle} \\
\midrule
Naissances & Injection $K_0 \to c\,K_0$ par construction. Le nombre de
naissances, tiré en $\mathrm{Poisson}(\lambda)$, est inchangé. \\
Choc & Multiplicatif : $cK \cdot e^{\xi}$, même $\xi$, donc $\times c$. \\
Production & $A'(cK)^{\gamma} = A' c^{\gamma} K^{\gamma}$, et
\eqref{eq:facteur} donne $c^{1-\gamma} = A'/A$, donc $A'c^{\gamma} = cA$ :
la production est $\times c$. \emph{C'est ici, et seulement ici, que
l'exposant $1/(1-\gamma)$ est requis.} \\
Taux d'intérêt & La règle marginale évalue $\gamma A K^{\gamma-1}$, qui
devient $\gamma A' c^{\gamma-1} K^{\gamma-1} = \gamma (A'c^{\gamma}/c)
K^{\gamma-1} = \gamma A K^{\gamma-1}$ : \textbf{invariant}. Un taux est un
nombre sans dimension. \\
Service & Principal $\times$ taux $= (cq)\,r$, donc $\times c$. \\
Dépréciation & Linéaire : $\times c$. \\
Institution & $\delta^{*} = h(C)-K$ avec $h$ homogène de degré 1 à
technologies fixées, donc $\times c$. \\
Faillites & Comparaisons entre grandeurs toutes en $c$ : les mêmes entités
meurent, aux mêmes pas. \\
\bottomrule
\end{tabularx}
\end{center}

\noindent Aucune phase ne consomme le générateur différemment --- naissances,
chocs et appariements du marché tirent le même nombre de valeurs, dans le
même ordre. Les deux simulations sont donc la même, à l'échelle $c$ près, et
\[
\mathrm{prod\_tot}(A') = c \cdot \mathrm{prod\_tot}(A)
\quad\Longrightarrow\quad
\varepsilon = \frac{\mathrm{d}\ln \mathrm{prod\_tot}}{\mathrm{d}\ln A}
= \frac{\ln c}{\ln (A'/A)} = \frac{1}{1-\gamma}.
\]

\paragraph{Ce qui empêche l'exactitude parfaite.} Trois constantes du moteur
sont dimensionnées et ne sont pas rééchelonnées : \code{MIN\_LOAN}
($10^{-9}$, transfert jugé négligeable), \code{ZERO\_TOL} (capital jugé nul)
et le plafond de transfert. Une paire dont le transfert optimal tombe entre
$\texttt{MIN\_LOAN}$ et $c\cdot\texttt{MIN\_LOAN}$ est refusée d'un côté et
acceptée de l'autre : c'est la seule origine possible d'un écart.

\paragraph{La vérification.} On lance depuis $t = 0$, à graine identique,
quatre variantes par $\gamma$ : la référence ($A=1$, $K_0=25$), la
compensation covariante ($A=1{,}5$, $K_0 = 25\,c$), l'absence de
compensation ($K_0$ inchangé), et une compensation au mauvais exposant
($K_0 \times 1{,}5$). On compare chaque trajectoire à celle de la référence
multipliée par $c$.

\fait{} La compensation covariante est exacte \textbf{à la précision
machine} : l'écart maximal sur 800 pas vaut $5{,}2\cdot10^{-15}$ sur
\code{prod\_tot} et $1{,}3\cdot10^{-15}$ sur \code{K\_tot} à
$\gamma = 0{,}5$, et la population est \emph{rigoureusement identique}
(écart exactement nul) --- les mêmes entités meurent aux mêmes pas, comme
l'argument le prédit. L'élasticité mesurée vaut $2{,}0000$ et $1{,}6667$,
c'est-à-dire $1/(1-\gamma)$ à toutes les décimales. \fait{} Les deux autres
compensations s'écartent de $24$ à $58\,\%$, et donnent
$\varepsilon = 0{,}71$ (aucune compensation) et $1{,}34$ (mauvais exposant).
\fait{} Le compteur \code{mkt\_blocked\_tiny} \emph{diffère} bien entre
référence et bras covariant (8 contre 7 à $\gamma=0{,}5$) : les seuils
dimensionnés mordent effectivement, mais sur des transferts inférieurs à
$10^{-9}$ joule, d'où le résidu à $10^{-15}$ et non à zéro.

\input{tables/scaling_theory.tex}

\incertitude{} \textbf{Cette exactitude ne s'étend pas au protocole de la
sous-section suivante}, et il ne faut pas confondre les deux. La covariance
vaut pour deux simulations lancées \emph{depuis $t=0$} avec des paramètres
homothétiques. La campagne §\ref{sec:loi} fait autre chose : elle branche
l'intervention à $t_0$ sur un état qui, lui, est encore à l'ancienne
échelle, et mesure la réponse pendant que le système s'y adapte. Les
$\pm1{,}3\,\%$ résiduels qu'on y trouve ne viennent donc pas des constantes
dimensionnées --- celles-ci pèsent $10^{-15}$ --- mais de la relaxation
incomplète et du nombre fini de graines. Leurs signes alternés
($-1{,}2$\,\%, $+0{,}9$\,\%, $-0{,}1$\,\%) sont cohérents avec du bruit et
non avec un biais de seuil.

\begin{figure}[H]
\centering
\includegraphics[width=\textwidth]{figures/scaling_covariance.png}
\caption{Vérification directe de la covariance. À gauche, $\gamma = 0{,}5$ :
la courbe de référence multipliée par $c$ (noir) et les trois compensations.
Seule $K_0 \times c$ se superpose. À droite, l'écart maximal à la
proportionnalité exacte pour les trois compensations et les trois $\gamma$,
en échelle logarithmique. Données :
\code{results/analysis/scaling\_theory.csv}.}
\label{fig:covariance}
\end{figure}

\subsection{Mise à l'épreuve sur trois valeurs de $\gamma$}

Une égalité vérifiée en un point n'est pas une loi. L'ablation donne
$\varepsilon = 2{,}018$ pour $\gamma = 0{,}5$, où $1/(1-\gamma)$ vaut
exactement 2 ; on rejoue donc le même protocole --- $A \times 1{,}5$ sur
toutes les entités, capital de naissance compensé par $A^{1/(1-\gamma)}$ ---
en $\gamma = 0{,}4$ et $\gamma = 0{,}6$, où la loi prédit $1{,}667$ et
$2{,}500$. Chaque $\gamma$ reçoit son propre amorçage de 2000 pas et son
propre contrôle apparié ; aucune intervention ne consomme de tirage, donc
l'appariement est exact.

\fait{} Contrôle préalable : l'amorçage est bien stationnaire à $t_0$ pour
ces deux nouveaux $\gamma$ --- le rapport entre la moyenne du dernier quart
de l'amorçage et celle du quart précédent vaut $1{,}0014$ et $0{,}9993$.
Le $t_0$ calibré sur $\gamma = 0{,}5$ y est donc valide.

\input{tables/scaling_gamma.tex}

\begin{figure}[H]
\centering
\includegraphics[width=\textwidth]{figures/scaling_gamma.png}
\caption{La loi d'échelle sur trois valeurs de $\gamma$. À gauche, les
élasticités mesurées contre la courbe $1/(1-\gamma)$ ; à droite, l'écart
relatif à la loi. Données : \code{results/analysis/scaling\_gamma.json}.}
\label{fig:loi}
\end{figure}

\fait{} La loi est vérifiée à $-1{,}2$~\%, $+0{,}9$~\% et $-0{,}1$~\% près sur
une plage où $\varepsilon$ lui-même varie de $1{,}65$ à $2{,}50$ --- soit une
variation de 52~\% de la quantité prédite, capturée à mieux que 1,3~\%.
\fait{} Et à chaque $\gamma$, la population reste celle de son contrôle
($\times0{,}993$, $\times1{,}005$, $\times0{,}999$) : la compensation
neutralise l'effet de population partout, pas seulement en $\gamma = 0{,}5$.

\inference{} L'énoncé qui résiste à ce test est donc : \textbf{lorsque le
capital de naissance suit l'échelle technologique, la production agrégée du
système se comporte comme son échelle autarcique}, $A^{1/(1-\gamma)}$ --- une
loi de puissance dont l'exposant est fixé par le seul rendement d'échelle de
la fonction de production, et non par les paramètres du marché du crédit.
\incertitude{} Testée ici sur $\gamma \in \{0{,}4;\,0{,}5;\,0{,}6\}$, à une
seule amplitude ($A \times 1{,}5$) et un seul jeu de paramètres de
population. Rien n'établit qu'elle survive à un $\gamma$ proche de 0 ou de 1,
ni à une amplitude très différente.

\subsection{Covariance de pas de temps : $\delta$, $\sigma$ et $\lambda$ ne
sont pas des paramètres de forme}
\label{sec:temps}

La covariance d'échelle du §\ref{sec:echelle} dit ce qui se passe quand on
change l'unité de \emph{capital}. Il en existe une seconde, qui dit ce qui se
passe quand on change l'unité de \emph{temps}, et qui a été proposée par le
relecteur sous cette forme : \emph{si l'on double $\lambda$, que l'on remplace
$\delta$ par $2\delta - \delta^2$ et $\sigma$ par $\sqrt{2}\,\sigma$, alors un
pas neuf vaut deux pas anciens} --- la population resterait la même, la
production serait grossièrement doublée.

\paragraph{Trois substitutions qui composent exactement.} \fait{} Les trois
formules proposées sont exactes, et chacune pour une raison différente.

\begin{description}
\item[Dépréciation.] Deux pas font $K(1-\delta)^2$ ; un pas neuf fait
  $K(1-\delta')$. L'égalité $1 - (2\delta - \delta^2) = (1-\delta)^2$ est une
  identité algébrique : \emph{aucune approximation}.
\item[Choc.] Le moteur tire $\xi \sim \mathcal{N}(-\sigma^2/2, \sigma)$ puis
  applique $K \leftarrow K e^{\xi}$ (\code{model.py:880--884}). Deux pas
  composent $e^{\xi_1 + \xi_2}$ avec
  $\xi_1 + \xi_2 \sim \mathcal{N}(-\sigma^2, \sigma\sqrt{2})$ ; un pas neuf
  avec $\sigma' = \sqrt{2}\sigma$ tire
  $\mathcal{N}(-\sigma'^2/2, \sigma') = \mathcal{N}(-\sigma^2, \sqrt{2}\sigma)$.
  Les deux lois coïncident, \emph{y compris la correction de dérive} --- qui
  est précisément la partie qu'un facteur $\sqrt{2}$ naïf raterait.
\item[Naissances.] Le moteur tire $\text{Poisson}(\lambda)$
  (\code{model.py:864}). Deux pas donnent
  $\text{Poisson}(\lambda) + \text{Poisson}(\lambda) = \text{Poisson}(2\lambda)$,
  un pas neuf $\text{Poisson}(2\lambda)$ : exact en loi.
\end{description}

\paragraph{Deux flux par pas que l'énoncé ne mentionne pas.} \fait{} Deux
phases du pas sont des \emph{débits}, et ne figurent pas dans la liste. La
\textbf{production} : deux pas anciens produisent $\approx 2AK^{\gamma}$, un
pas neuf n'en produit que $AK^{\gamma}$ --- il manque un facteur 2 sur $A$.
Le \textbf{marché} : le moteur fait $R = \lfloor \rho N \rfloor$ rondes
d'appariement par pas (\code{model.py:497}), donc deux pas font $2\rho N$
rondes et un pas neuf $\rho N$ --- il manque un facteur 2 sur $\rho$.
\inference{} Et $A$ porte les deux flux monétaires à la fois : le taux
d'intérêt vaut $\gamma A K^{\gamma-1}$, donc doubler $A$ double aussi le
service par pas, ce qu'exige la composition.

\fait{} La conséquence se calcule avant de simuler. Le point fixe autarcique
passe de $[A(1-\delta)/\delta]^{1/(1-\gamma)}$ à
$[A'(1-\delta)^2/(2\delta - \delta^2)]^{1/(1-\gamma)}$, soit un facteur
\[
\left[\frac{A'}{A}\cdot\frac{1-\delta}{2-\delta}\right]^{\frac{1}{1-\gamma}}
\;=\;
\begin{cases}
0{,}2475 & \text{sans } A' = 2A ,\\[2pt]
0{,}9900 & \text{avec } A' = 2A .
\end{cases}
\]
\inference{} Le recalage n'est donc \textbf{jamais exact}. Ce déplacement-là
est d'ordre $\delta$ et s'annulerait avec lui --- mais il ne vaut que
$1\,\%$, et on verra plus bas que le résidu mesuré est dix fois plus grand :
$\Kaut$ n'est pas la source principale de l'erreur. Sans le facteur sur $A$,
en revanche, il n'y a plus de petit paramètre du tout : l'échelle autarcique
est divisée par 4.

\paragraph{Ce que la mesure donne.} Quatre bras lancés depuis $t = 0$, trois
graines, comparés à \emph{temps ancien égal} --- le pas $t$ d'un bras rapide
répond au pas $2t$ de la référence, et les flux des bras rapides sont divisés
par 2 pour être rapportés au pas ancien. \fait{} Contrôle préalable : les
quatre bras sont stationnaires à l'arrivée, le rapport du dernier quart au
quart précédent valant $0{,}993$ à $1{,}003$ sur \code{K\_tot},
\code{pop} et \code{prod\_tot} --- les niveaux ci-dessous sont des niveaux et
non des transitoires.

\input{tables/time_rescaling.tex}

\noindent\footnotesize Chaque cellule est le rapport de l'observable à celle
du bras \code{ref} à temps ancien égal, moyenné sur la fenêtre des 800
derniers pas anciens ; le $\pm$ est l'écart-type sur les trois graines. Un
recalage parfait donnerait $1$ partout. Dernière ligne : le déplacement de
$\Kaut$ calculé analytiquement, sans simuler.\normalsize

\begin{figure}[H]
\centering
\includegraphics[width=\textwidth]{figures/time_rescaling.png}
\caption{Recalage temporel, graine 0. À gauche et au centre, les niveaux en
temps ancien ; à droite, le rapport à la référence, qu'une échelle
logarithmique rendrait invisible. Données :
\code{results/analysis/time\_rescaling.json}.}
\label{fig:temps}
\end{figure}

\fait{} \textbf{L'énoncé brut ne conserve pas le système.} La population monte
à $\times 1{,}853$, la production par pas ancien tombe à $\times 0{,}701$, la
durée de vie s'allonge à $\times 1{,}859$ et la tension passe de $13{,}3$ à
$5{,}7$ ($\times 0{,}434$). \inference{} Le capital total, lui, ne bouge
presque pas ($\times 1{,}058$) --- non parce que le système est conservé, mais
parce que les deux effets se compensent : chaque entité est deux fois plus
petite et il y en a presque deux fois plus. C'est un piège de lecture
qu'il vaut la peine de nommer : \emph{sur cette transformation, \code{K\_tot}
est l'observable la moins informative de toutes}.

\fait{} \textbf{Avec les deux flux, le recalage tient.} Population
$\times 1{,}067$, morts par pas ancien $\times 1{,}003$, durée de vie
$\times 1{,}064$, tension $\times 0{,}943$, production par pas ancien
$\times 1{,}118$ --- soit, par pas \emph{neuf}, une production multipliée par
$2{,}24$. \inference{} Les trois prédictions du relecteur sont donc vérifiées
dans leur substance : population inchangée à $7\,\%$ près, production
« grossièrement multipliée par deux », changements marginaux. Elles ne le sont
pas dans leur lettre : sans $A \times 2$ ni $\rho \times 2$, aucune des trois
ne tient.

\fait{} Le rôle de $\rho$ se lit sur le bras intermédiaire : avec
$A \times 2$ mais $\rho$ inchangé, la population est déjà rattrapée
($\times 1{,}145$) mais le nombre de prêts actifs tombe à $\times 0{,}690$ ---
le marché du crédit tourne deux fois moins vite que le reste du système.
Le rétablir ramène ce nombre à $\times 1{,}167$.

\incertitude{} \textbf{Ce qui reste, et d'où ça vient.} Le résidu de $6$ à
$17\,\%$ du bras complet n'est pas du bruit : les $\pm$ du tableau sont des
écarts-types sur trois graines, et ils valent $0{,}4$ à $1{,}4\,\%$ sur ce
bras. Il a deux sources, et \textbf{ce n'est pas la première qui domine}.

\fait{} La discrétisation calculée plus haut déplace $\Kaut$ de $1{,}0\,\%$.
Le résidu observé vaut $11{,}8\,\%$ sur la production, $15{,}8\,\%$ sur le
capital et $16{,}7\,\%$ sur le nombre de prêts : \textbf{un ordre de grandeur
au-dessus}. \inference{} La discrétisation en $\delta$ n'explique donc au
mieux qu'un dixième de l'écart. Le reste vient de la seconde source, la
\textbf{phase de marché, qui ne se compose pas} : deux rondes de $\rho N$
appariements ne valent pas une ronde de $2\rho N$, parce que les appariements
de la seconde voient l'état laissé par la première --- des prêts déjà noués,
des capitaux déjà déplacés. Aucun choix de paramètres ne répare cela.

\paragraph{Deux tests, dont un qui ne tranche pas.} La transformation se
généralise sans difficulté à un \emph{pas de recalage} $s$ quelconque --- le
nombre de pas anciens comprimés en un :
\[
\delta_s = 1 - (1-\delta)^s, \quad \sigma_s = \sigma\sqrt{s}, \quad
\lambda_s = s\lambda, \quad A_s = sA, \quad \rho_s = s\rho .
\]

\fait{} \textbf{Premier test : $s = 4$.} Le résidu triple --- production
$+11{,}8\,\%$ puis $+34{,}6\,\%$, capital $+15{,}8\,\%$ puis $+48{,}0\,\%$,
prêts actifs $+16{,}7\,\%$ puis $+53{,}9\,\%$. \incertitude{} \textbf{Cela ne
prouve rien}, et il faut le dire : le déplacement de $\Kaut$ triple lui aussi
($1{,}0\,\%$ puis $3{,}0\,\%$), parce que les deux termes croissent en
$(s-1)$ --- un pas neuf recouvre $s-1$ compositions manquées, qu'il s'agisse
de dépréciation ou d'appariement. Le test est compatible avec les deux
explications. Ce qu'il montre seulement, c'est que le rapport reste stable :
le terme de $\Kaut$ pèse $8\,\%$ du résidu à $s = 2$ comme à $s = 4$.

\fait{} \textbf{Second test : $\delta$ divisé par 5, à $s = 2$ inchangé.}
Celui-là sépare, parce que le terme de discrétisation est proportionnel à
$\delta$ tandis que le marché ne le voit pas. À $\delta = 0{,}002$, $\Kaut$
ne se déplace plus que de $0{,}20\,\%$ au lieu de $1{,}00\,\%$ --- exactement
cinq fois moins. Si le résidu était de discrétisation, il devrait tomber d'un
facteur 5 ; s'il vient du marché, il ne doit pas bouger.

\input{tables/time_rescaling_stride.tex}

\fait{} \textbf{Le résidu ne bouge pas.} À $\delta$ cinq fois plus petit, la
production passe de $+11{,}8\,\%$ à $+15{,}0\,\%$, le capital de
$+15{,}8\,\%$ à $+22{,}1\,\%$, les prêts actifs de $+16{,}7\,\%$ à
$+17{,}5\,\%$ --- aucune décroissance, et même une légère hausse, quand la
prédiction « discrétisation » demandait une division par 5.
\inference{} La discrétisation en $\delta$ est donc écartée comme source
principale : \textbf{le résidu vient de la phase de marché}. Ce n'est plus
une explication par élimination, c'est une prédiction quantitative qui a été
mise à l'épreuve et qui a survécu au bon test après avoir échappé au mauvais.

\inference{} C'est exactement ce qui distingue ce modèle d'une équation
différentielle, et c'est un résultat en soi : \textbf{le pas de temps n'est
pas qu'une unité, c'est aussi le grain auquel le marché se réapparie}. Un
recalage temporel exact demanderait de composer les appariements eux-mêmes,
c'est-à-dire de savoir ce que $s$ rondes successives font qu'une ronde $s$
fois plus grosse ne fait pas --- question ouverte, et probablement le vrai
contenu de la limite en temps continu de ce modèle.

\inference{} \textbf{Ce que l'on retient pour la suite.} Les paramètres du
modèle se rangent en deux classes. $\delta$, $\sigma$, $\lambda$, $A$ et
$\rho$ sont des grandeurs \emph{par pas} : elles se recalent avec l'unité de
temps, et une part de leur effet apparent n'est qu'un changement d'unité.
$\gamma$, $K_0$, la règle de transfert et la règle de taux sont des grandeurs
\emph{de forme} : elles n'en dépendent pas. Rapproché du
§\ref{sec:echelle}, cela fait deux groupes de covariance --- un sur le
capital, un sur le temps --- qui réduisent d'autant le nombre de paramètres
réellement libres du modèle.

\section{Verdict}

\begin{center}
\begin{tabularx}{\textwidth}{@{}lXX@{}}
\toprule
& \textbf{court horizon} ($h \lesssim 300$) & \textbf{régime établi} ($h \ge 1001$) \\
\midrule
\textbf{portée partielle} (\code{fraction}) &
\textbf{Effet rebond OBSERVÉ.} Réponse 2,3 à 2,9 fois la réponse
proportionnelle, 3 à 48$\sigma$, stable sur deux amplitudes, trois
intensités et deux leviers. &
\textbf{NON TRANCHABLE.} La cohorte traitée s'éteint (part de production
$0{,}198 \to 0{,}022$) ; les écarts restent sous $2\sigma$ et le
dénominateur de l'élasticité disparaît. \\
\addlinespace
\textbf{portée globale} (\code{all}, \code{new}), $K_0$ fixe &
\textbf{Rebond transitoire observé} : $E = 2{,}18$ à $h\le50$
($+108{,}8$\,\% pour un choc de $+50$\,\%). &
\textbf{Effet INVERSE.} $E = 0{,}72$--$0{,}73$ ($88\sigma$) : $+36{,}1$\,\%
de production pour $+50$\,\% de $A$. Par entité $\times1{,}80$, mais
population $\times0{,}754$. \\
\addlinespace
\textbf{portée globale}, $K_0$ compensé à l'échelle (§\ref{sec:ablation}) &
Rebond transitoire également. &
\textbf{Rebond ÉTABLI, fort.} $E = 2{,}53$, $\varepsilon = 2{,}02 =
1/(1-\gamma)$ : $+126{,}7$\,\% de production pour $+50$\,\% de $A$,
à population inchangée. \\
\bottomrule
\end{tabularx}
\end{center}

\noindent Trois lectures s'imposent, dans cet ordre :

\begin{enumerate}
\item \textbf{Le rebond existe dans ce modèle, mais c'est un phénomène de
  transitoire.} Toutes les portées et toutes les amplitudes produisent la
  même amplification de facteur $\approx$2,5 dans les dizaines de pas qui
  suivent l'intervention.
\item \textbf{En régime établi à dotation de naissance inchangée, une
  hausse globale de la puissance d'extraction rend le système plus petit.}
  L'agrégat capte moins que proportionnellement la hausse, non parce que les
  entités produisent moins, mais parce qu'il y en a un quart de moins.
\item \textbf{Cette contraction n'est pas un effet de la technologie : c'est
  un effet d'échelle du capital de naissance.} L'ablation du
  §\ref{sec:ablation} l'annule entièrement en faisant suivre $K_0$, et la
  réponse devient alors fortement super-proportionnelle. Le signe même du
  résultat global dépend donc d'une convention --- ce que le prompt laissait
  ouvert et que cette itération a tranché par la mesure.
\item \textbf{Une intervention partielle ne laisse pas de trace
  permanente.} Sans mécanisme de transmission aux naissances, la cohorte
  dopée est absorbée par le renouvellement en quelques centaines de pas.
\item \textbf{Une grandeur résume ces effets d'échelle, et une seule
  hypothèse la met en défaut.} La tension $T = \Kaut/K_{\text{eq}}$
  (§\ref{sec:tension}) rend la mortalité prévisible à $1{,}4\,\%$ près sur
  une plage de tension $\times15$ obtenue par trois leviers différents ---
  \emph{à condition que la dépréciation soit fixe}. Le levier $\delta$ la met
  en échec sans ambiguïté (§\ref{sec:tension_test}). Ce que les quatre
  leviers alignent, en revanche, c'est la rotation du crédit, avec
  $R^2 = 0{,}998$. \incertitude{} C'est une variable endogène : elle désigne
  le canal, elle ne le démontre pas.
\item \textbf{L'exposant $1/(1-\gamma)$ n'est pas un ajustement.} Il découle
  d'une covariance d'échelle exacte du modèle, démontrée au
  §\ref{sec:echelle} et vérifiée numériquement à la précision machine. Les
  écarts résiduels de la loi empirique s'expliquent par les seules
  constantes dimensionnées non rééchelonnées du moteur.
\end{enumerate}

\section{Limites}

\begin{itemize}
\item \textbf{Observation, pas attribution.} Conformément au §1 du prompt,
  aucune décomposition causale n'est revendiquée. Les canaux du §3.5 sont
  des mesures concomitantes.
\item \textbf{Un seul jeu de paramètres de population.} $\lambda=30$,
  $\delta=0{,}01$, $\sigma=0{,}01$, $K_0=25$, $\gamma=0{,}5$, $A=1$. Le §2
  du prompt suggérait d'explorer d'autres régimes de population/marché en
  axe secondaire ; cela n'a pas été fait et reste ouvert.
\item \textbf{Cinq graines, et ce que cela coûte.} Suffisant pour les effets
  à 20--88$\sigma$, juste suffisant pour distinguer les fenêtres longues du
  bruit. Avec $n=5$, la statistique $\bar X/\mathrm{SE}$ suit une Student à
  4 degrés de liberté et non une normale (note~\ref{note:sigma}) : le seuil
  à 5\,\% est $2{,}78$ et non $1{,}96$. Aucune conclusion de ce rapport ne
  bascule à cette lecture, mais la marge est mince sur la cellule
  $\varphi=0{,}05$ ($3{,}2\sigma$, $p=3{,}3\,\%$). Douze graines
  ramèneraient le seuil à $2{,}20$ pour un coût de calcul multiplié par
  $2{,}4$.
\item \textbf{$K_0$ délibérément non co-varié.} Une hausse de $A$ déplace
  l'échelle autarcique ($\Kaut \propto A^{1/(1-\gamma)}$) : les bras
  globaux changent donc simultanément l'extraction et la tension du
  système par rapport à cette échelle. La décision et son motif sont
  documentés dans le rapport de conception ; le caveat demeure, et le
  §\ref{sec:ablation} en mesure la conséquence.
\item \textbf{La tension est mesurée, pas instrumentée.} $K_{\text{eq}}$ est
  reconstitué depuis les agrégats de \code{tech\_series}, dont l'effectif
  est celui de \emph{fin} de pas ; l'effectif producteur est reconstruit
  (colonne \code{basis}), exactement quand une seule technologie est
  vivante, au prorata sinon. Enregistrer l'effectif et le capital
  producteurs directement dans le moteur rendrait la mesure exacte en toute
  circonstance ; cela demande une modification du moteur, écartée ici.
\item \textbf{Amplitude dérivante du bras $\gamma$.} $K^{\Delta\gamma}$
  croît avec le capital ; l'élasticité de ce bras à long horizon est la
  plus fragile du tableau.
\item \textbf{Une seule institution testée en campagne.}
  \code{transfer\_cap="optimum"} et \code{rate\_rule="marginal"}. Les deux
  variantes implémentées (plafond à l'égalisation, taux au partage du
  surplus) n'ont été piloté-es que ponctuellement.
\item \textbf{Le régime (c) n'est exercé que par un bras.} Le chemin
  coûteux du noyau (exposants différents, table 1D) est bien sollicité,
  mais par une seule cellule.
\end{itemize}

\section{Reproduire}

\begin{quote}\footnotesize
\begin{verbatim}
cd /home/anatole/jupyter/m4_3live_credit_soc
python3 scripts/campaign.py burn     # 5 amorçages 0 -> t0=2000  (~220 s, 5 procs)
python3 scripts/campaign.py arms     # 9 bras x 5 graines        (~1865 s, 6 procs)
python3 scripts/ablation_k0.py run   # 5 bras x 5 graines        (~1287 s, 6 procs)
python3 scripts/scaling_gamma.py run # loi d'echelle, 2 gammas
python3 scripts/tension_sweep.py run # balayage a trois leviers  (~3900 s, 6 procs)
python3 scripts/scaling_theory.py    # covariance d'echelle      (~600 s, 6 procs)
python3 scripts/tension_vs_A.py run  # balayage en A, 2 gammas   (~1200 s, 6 procs)
python3 scripts/time_rescaling.py    # covariance de pas de temps (~390 s, 6 procs)

python3 scripts/analyse.py           # metriques, figures, controles d'integrite
python3 scripts/exact_amplitude.py   # amplitude exacte, entite par entite (~15 s)
python3 scripts/tension_figures.py   # tension.csv + figure, tous les runs
python3 scripts/tension_analysis.py  # lecture croisee de la tension
python3 scripts/tension_vs_A.py      # exposant dlnT/dlnA, par gamma
python3 scripts/verification_figures.py
python3 scripts/make_report_tables.py
python3 scripts/import_to_simulation_lab.py
python3 scripts/make_traceability.py
\end{verbatim}
\end{quote}

Données brutes : \code{results/campaign/}. Métriques agrégées :
\code{results/analysis/metrics.json} et \code{windows.csv}. Chaque bras
conserve son \code{interventions.jsonl}, rejouable à l'identique par
\code{driver/headless.py replay}.

\end{document}
